\PassOptionsToPackage{no-math}{fontspec}
% texlate: native XeTeX font and PDF-driver capabilities
\AddToHook{package/microtype/after}{%
\DeclareMicrotypeSet{texlate-native}{encoding={TU,EU1,EU2}}%
\UseMicrotypeSet[protrusion]{texlate-native}%
}
% breakurl already has a native-PDF branch; limit its PDF-mode override to loading.
\AddToHook{package/breakurl/before}{%
\RequirePackage{xkeyval,ifpdf}%
\let\TeXlateSavedIfpdf\ifpdf\let\ifpdf\iftrue
}
\AddToHook{package/breakurl/after}{\let\ifpdf\TeXlateSavedIfpdf}
% This class's optional arXiv check assumes every non-pdfTeX engine writes DVI.
\PassOptionsToClass{nopdfoutputerror,allowfontchangeintitle}{quantumarticle}
% Embedded PostScript can silently disappear with restricted XeTeX drivers.
% Record actual drawing operations; merely loading an unused package is harmless.
\AddToHook{package/pstricks/after}{%
\ifcsname pst@object\endcsname
\expandafter\let\expandafter\TeXlatePstObject\csname pst@object\endcsname
\expandafter\def\csname pst@object\endcsname#1{%
\typeout{TeXlate-PostScript-object: #1}\TeXlatePstObject{#1}}%
\fi
}
% \DeclareUnicodeCharacter is pdfTeX/inputenc-only and undefined under XeTeX,
% but e-print preambles still call it (2501.14787). Emulate via the lccode
% idiom: make the code point an active char expanding to the replacement.
\providecommand{\DeclareUnicodeCharacter}[2]{%
\begingroup\lccode`\~="#1\relax
\lowercase{\endgroup\catcode`~\active\protected\def~}{#2}}
\documentclass{acm_proc_article-sp}
\usepackage[fontset=fandol,UTF8]{ctex}  % [texlate injected]
% texlate: twin named-dest for shared-counter theorems
\makeatletter
\def\TeXlate@thmtwin{%
  \@ifundefined{MakeLinkTarget}{}{%
  \@ifundefined{@currentcounter}{}{%
  \@ifundefined{theH\@currentcounter}{}{%
    \ifx\@currenvir\@currentcounter
      \@ifundefined{alias@ctr@\@currenvir}{}{%
        \edef\TeXlate@twin{\csname alias@ctr@\@currenvir\endcsname}%
        \TeXlate@emit}%
    \else
      \@ifundefined{the\@currenvir}{}{%
        \let\TeXlate@twin\@currenvir
        \TeXlate@emit}%
    \fi}}}%
}%
\def\TeXlate@emit{%
  \begingroup
  \let\TeXlate@save\@currentHref
  \edef\TeXlate@name{\TeXlate@twin.\csname theH\@currentcounter\endcsname}%
  \MakeLinkTarget*{\TeXlate@name}%
  \global\let\@currentHref\TeXlate@save
  \endgroup}%
\AddToHook{cmd/@begintheorem/before}{\TeXlate@thmtwin}%
\AddToHook{cmd/@opargbegintheorem/before}{\TeXlate@thmtwin}%
\makeatother

% texlate: math fallback via dedicated symbol fonts
\ifdefined\Umathcode
\makeatletter
\DeclareFontFamily{TU}{texlatecjk}{\hyphenchar\font\m@ne}
\DeclareFontShape{TU}{texlatecjk}{m}{n}{<->"[FandolSong-Regular.otf]"}{}
\DeclareFontShape{TU}{texlatecjk}{b}{n}{<->"[FandolSong-Bold.otf]"}{}
\DeclareFontShape{TU}{texlatecjk}{bx}{n}{<->ssub*texlatecjk/b/n}{}
\DeclareFontShape{TU}{texlatecjk}{m}{it}{<->ssub*texlatecjk/m/n}{}
\DeclareFontShape{TU}{texlatecjk}{m}{sl}{<->ssub*texlatecjk/m/n}{}
\DeclareFontShape{TU}{texlatecjk}{b}{it}{<->ssub*texlatecjk/b/n}{}
\DeclareFontShape{TU}{texlatecjk}{bx}{it}{<->ssub*texlatecjk/b/n}{}
\DeclareFontShape{TU}{texlatecjk}{bx}{sl}{<->ssub*texlatecjk/b/n}{}
\DeclareSymbolFont{texlatecjk}{TU}{texlatecjk}{m}{n}
\SetSymbolFont{texlatecjk}{bold}{TU}{texlatecjk}{b}{n}
\def\TeXlate@mathmap#1#2-#3;{\count@="#2\relax
  \@whilenum\count@<"#3 \do{\Umathcode\count@="0 #1 \count@\advance\count@\@ne}}
\TeXlate@mathmap\symtexlatecjk 4E00-9FFF;
\TeXlate@mathmap\symtexlatecjk 3400-4DBF;
\TeXlate@mathmap\symtexlatecjk 3000-303F;
\TeXlate@mathmap\symtexlatecjk FF00-FFEF;
\TeXlate@mathmap\symtexlatecjk 3040-30FF;
\TeXlate@mathmap\symtexlatecjk F900-FAFF;
\TeXlate@mathmap\symtexlatecjk 2E80-2FDF;
\TeXlate@mathmap\symtexlatecjk 20000-2A6DF;
\TeXlate@mathmap\symtexlatecjk 2A700-2EBEF;
% 非 CJK 带（西里尔人名 Ш/Д/Л、组合符、拉丁扩展）走 Libertinus Serif;
% \IfFileExists 门——otf 缺席则整段不挂, 免把缺字升级成字体加载错误。
% 00D7 ×/00F7 ÷ 是 binop 语义, 普通 ordinary 化会改距, 故带内挖掉。
\IfFileExists{LibertinusSerif-Regular.otf}{%
\DeclareFontFamily{TU}{texlatefb}{\hyphenchar\font\m@ne}
\DeclareFontShape{TU}{texlatefb}{m}{n}{<->"[LibertinusSerif-Regular.otf]"}{}
\DeclareFontShape{TU}{texlatefb}{b}{n}{<->"[LibertinusSerif-Bold.otf]"}{}
\DeclareFontShape{TU}{texlatefb}{bx}{n}{<->ssub*texlatefb/b/n}{}
\DeclareFontShape{TU}{texlatefb}{m}{it}{<->ssub*texlatefb/m/n}{}
\DeclareFontShape{TU}{texlatefb}{m}{sl}{<->ssub*texlatefb/m/n}{}
\DeclareFontShape{TU}{texlatefb}{b}{it}{<->ssub*texlatefb/b/n}{}
\DeclareFontShape{TU}{texlatefb}{bx}{it}{<->ssub*texlatefb/b/n}{}
\DeclareFontShape{TU}{texlatefb}{bx}{sl}{<->ssub*texlatefb/b/n}{}
\DeclareSymbolFont{texlatefb}{TU}{texlatefb}{m}{n}
\SetSymbolFont{texlatefb}{bold}{TU}{texlatefb}{b}{n}
\TeXlate@mathmap\symtexlatefb 0400-04FF;
\TeXlate@mathmap\symtexlatefb 0300-036F;
\TeXlate@mathmap\symtexlatefb 00C0-00D6;
\TeXlate@mathmap\symtexlatefb 00D8-00F6;
\TeXlate@mathmap\symtexlatefb 00F8-017F;
}{}
\makeatother
\fi

% texlate: xeCJK first-use warmup (frontmatter \vbox poison)
\ifdefined\XeTeXversion
\AtBeginDocument{\setbox0=\hbox{字}}%
\fi

% texlate: \t absent from tuenc.def -> TS1 \accent bypasses xeCJK
\ifdefined\UnicodeEncodingName
\DeclareUnicodeAccent{\t}{"0361}
\fi

% \usepackage{refcheck}

%\usepackage{fullpage}
%\usepackage{times}
%\usepackage{microtype}
\usepackage{amsmath,amsfonts,amssymb}
\makeatletter\let\th@plain\@undefined\makeatother
\usepackage[amsmath,thmmarks,hyperref]{ntheorem}
\usepackage{mathtools}
%\usepackage{url}
\usepackage{enumitem}
%\usepackage{hyperref}
\usepackage{xspace}
%\usepackage{stmaryrd}
\usepackage{nag}
%\usepackage{tikz}
\usepackage{comment}


% bring in the individual header files

% blackboard symbols

\newcommand{\C}{\ensuremath{\mathbb{C}}}
\newcommand{\D}{\ensuremath{\mathbb{D}}}
\newcommand{\F}{\ensuremath{\mathbb{F}}}
\newcommand{\G}{\ensuremath{\mathbb{G}}}
%\newcommand{\J}{\ensuremath{\mathbb{J}}}
\newcommand{\N}{\ensuremath{\mathbb{N}}}
\newcommand{\Q}{\ensuremath{\mathbb{Q}}}
\newcommand{\R}{\ensuremath{\mathbb{R}}}
\newcommand{\T}{\ensuremath{\mathbb{T}}}
\newcommand{\Z}{\ensuremath{\mathbb{Z}}}

\newcommand{\Zt}{\ensuremath{\Z_t}}
\newcommand{\Zp}{\ensuremath{\Z_p}}
\newcommand{\Zq}{\ensuremath{\Z_q}}
\newcommand{\ZN}{\ensuremath{\Z_N}}
\newcommand{\Zps}{\ensuremath{\Z_p^*}}
\newcommand{\ZNs}{\ensuremath{\Z_N^*}}
\newcommand{\JN}{\ensuremath{\J_N}}
\newcommand{\QRN}{\ensuremath{\mathbb{QR}_N}}
\newcommand{\Fq}{\ensuremath{\F_q}}

% "left-right" pairs of symbols

%% NOTE: this requires \usepackage{mathtools} in the document preamble

% inner product
\DeclarePairedDelimiter\inner{\langle}{\rangle}
% absolute value
\DeclarePairedDelimiter\abs{\lvert}{\rvert}
% a set
\DeclarePairedDelimiter\set{\{}{\}}
% parens
\DeclarePairedDelimiter\parens{(}{)}
% tuple, alias for parens
\DeclarePairedDelimiter\tuple{(}{)}
% square brackets
\DeclarePairedDelimiter\bracks{[}{]}
% rounding off
\DeclarePairedDelimiter\round{\lfloor}{\rceil}
% floor function
\DeclarePairedDelimiter\floor{\lfloor}{\rfloor}
% ceiling function
\DeclarePairedDelimiter\ceil{\lceil}{\rceil}
% length of some vector, element
\DeclarePairedDelimiter\length{\lVert}{\rVert}
% "lifting" of a residue class
\DeclarePairedDelimiter\lift{\llbracket}{\rrbracket}

% macros for matrices and vectors
\newcommand{\matA}{\ensuremath{\mathbf{A}}}
\newcommand{\matB}{\ensuremath{\mathbf{B}}}
\newcommand{\matC}{\ensuremath{\mathbf{C}}}
\newcommand{\matD}{\ensuremath{\mathbf{D}}}
\newcommand{\matE}{\ensuremath{\mathbf{E}}}
\newcommand{\matF}{\ensuremath{\mathbf{F}}}
\newcommand{\matG}{\ensuremath{\mathbf{G}}}
\newcommand{\matH}{\ensuremath{\mathbf{H}}}
\newcommand{\matI}{\ensuremath{\mathbf{I}}}
\newcommand{\matJ}{\ensuremath{\mathbf{J}}}
\newcommand{\matK}{\ensuremath{\mathbf{K}}}
\newcommand{\matL}{\ensuremath{\mathbf{L}}}
\newcommand{\matM}{\ensuremath{\mathbf{M}}}
\newcommand{\matN}{\ensuremath{\mathbf{N}}}
\newcommand{\matP}{\ensuremath{\mathbf{P}}}
\newcommand{\matQ}{\ensuremath{\mathbf{Q}}}
\newcommand{\matR}{\ensuremath{\mathbf{R}}}
\newcommand{\matS}{\ensuremath{\mathbf{S}}}
\newcommand{\matT}{\ensuremath{\mathbf{T}}}
\newcommand{\matU}{\ensuremath{\mathbf{U}}}
\newcommand{\matV}{\ensuremath{\mathbf{V}}}
\newcommand{\matW}{\ensuremath{\mathbf{W}}}
\newcommand{\matX}{\ensuremath{\mathbf{X}}}
\newcommand{\matY}{\ensuremath{\mathbf{Y}}}
\newcommand{\matZ}{\ensuremath{\mathbf{Z}}}
\newcommand{\matzero}{\ensuremath{\mathbf{0}}}

\newcommand{\veca}{\ensuremath{\mathbf{a}}}
\newcommand{\vecb}{\ensuremath{\mathbf{b}}}
\newcommand{\vecc}{\ensuremath{\mathbf{c}}}
\newcommand{\vecd}{\ensuremath{\mathbf{d}}}
\newcommand{\vece}{\ensuremath{\mathbf{e}}}
\newcommand{\vecf}{\ensuremath{\mathbf{f}}}
\newcommand{\vecg}{\ensuremath{\mathbf{g}}}
\newcommand{\vech}{\ensuremath{\mathbf{h}}}
\newcommand{\veck}{\ensuremath{\mathbf{k}}}
\newcommand{\vecm}{\ensuremath{\mathbf{m}}}
\newcommand{\vecp}{\ensuremath{\mathbf{p}}}
\newcommand{\vecq}{\ensuremath{\mathbf{q}}}
\newcommand{\vecr}{\ensuremath{\mathbf{r}}}
\newcommand{\vecs}{\ensuremath{\mathbf{s}}}
\newcommand{\vect}{\ensuremath{\mathbf{t}}}
\newcommand{\vecu}{\ensuremath{\mathbf{u}}}
\newcommand{\vecv}{\ensuremath{\mathbf{v}}}
\newcommand{\vecw}{\ensuremath{\mathbf{w}}}
\newcommand{\vecx}{\ensuremath{\mathbf{x}}}
\newcommand{\vecy}{\ensuremath{\mathbf{y}}}
\newcommand{\vecz}{\ensuremath{\mathbf{z}}}
\newcommand{\veczero}{\ensuremath{\mathbf{0}}}

\theoremstyle{plain}            % following are "theorem" style
\newtheorem{theorem}{Theorem}[section]
\newtheorem{lemma}[theorem]{Lemma}
\newtheorem{corollary}[theorem]{Corollary}
\newtheorem{proposition}[theorem]{Proposition}
\newtheorem{claim}[theorem]{Claim}
\newtheorem{fact}[theorem]{Fact}
\newtheorem{techlemma}[theorem]{Technical Lemma}

\newtheorem{maintheorem}{Main Theorem}

\theoremstyle{definition}       % following are def style
\newtheorem{definition}[theorem]{Definition}
\newtheorem{conjecture}[theorem]{Conjecture}

\theoremstyle{remark}           % following are remark style
\newtheorem{remark}[theorem]{Remark}
\newtheorem{example}[theorem]{Example}
\newtheorem{note}[theorem]{Note}

% equation numbering style
\numberwithin{equation}{section}

% GENERAL COMPUTING STUFF

\newcommand{\bit}{\ensuremath{\set{0,1}}}
\newcommand{\pmone}{\ensuremath{\set{-1,1}}}

% asymptotic stuff
\DeclareMathOperator{\poly}{poly}
\DeclareMathOperator{\polylog}{polylog}
\DeclareMathOperator{\negl}{negl}
\newcommand{\Otil}{\ensuremath{\tilde{O}}}

% probability/distribution stuff
\DeclareMathOperator*{\E}{\mathbb{E}}
\DeclareMathOperator*{\Var}{Var}

% assorted
\DeclareMathOperator*{\wt}{wt}

% hash functions
\newcommand{\calH}{\ensuremath{\mathcal{H}}}
\newcommand{\calR}{\ensuremath{\mathcal{R}}}
\newcommand{\calX}{\ensuremath{\mathcal{X}}}
\newcommand{\calY}{\ensuremath{\mathcal{Y}}}

\newcommand{\compind}{\ensuremath{\stackrel{c}{\approx}}}
\newcommand{\statind}{\ensuremath{\stackrel{s}{\approx}}}
\newcommand{\perfind}{\ensuremath{\equiv}}

% font for general-purpose algorithms
\newcommand{\algo}[1]{\ensuremath{\mathsf{#1}}\xspace}

% font for general-purpose computational problems
\newcommand{\problem}[1]{\ensuremath{\mathsf{#1}}\xspace}

% font for complexity classes
\newcommand{\class}[1]{\ensuremath{\mathsf{#1}}\xspace}

% complexity classes and languages
\renewcommand{\P}{\class{P}}
\newcommand{\BPP}{\class{BPP}}
\newcommand{\NP}{\class{NP}}
\newcommand{\coNP}{\class{coNP}}
\newcommand{\AM}{\class{AM}}
\newcommand{\coAM}{\class{coAM}}
\newcommand{\IP}{\class{IP}}
\newcommand{\SZK}{\class{SZK}}
\newcommand{\NISZK}{\class{NISZK}}
\newcommand{\NICZK}{\class{NICZK}}
\newcommand{\PPP}{\class{PPP}}
\newcommand{\PPAD}{\class{PPAD}}

\newcommand{\yes}{\ensuremath{\text{YES}}}
\newcommand{\no}{\ensuremath{\text{NO}}}

\newcommand{\Piyes}{\ensuremath{\Pi^{\yes}}}
\newcommand{\Pino}{\ensuremath{\Pi^{\no}}}

% basic notation

\newcommand{\lamperp}{\Lambda^{\perp}}
\newcommand{\lam}{\Lambda}

% lattice
\newcommand{\lat}{\mathcal{L}}
% fundamental region
\newcommand{\piped}{\mathcal{P}}
% smoothing parameter
\newcommand{\smooth}{\eta}
% smoothing w/ epsilon
\newcommand{\smootheps}{\smooth_{\epsilon}}
% ball
\newcommand{\ball}{\mathcal{B}}
% cube
\newcommand{\cube}{\mathcal{C}}
% covering radius symbol
\newcommand{\cover}{\ensuremath{\mu}}
% Gram-Schmidt
\newcommand{\gs}[1]{\ensuremath{\widetilde{#1}}}
% GS minimum
\newcommand{\gsmin}{\ensuremath{\widetilde{bl}}}
% volume operation
\DeclareMathOperator{\vol}{vol}
% Hermite normal form
\DeclareMathOperator{\hnf}{HNF}
% rank
\DeclareMathOperator{\rank}{rank}
% distance operator
\DeclareMathOperator{\dist}{dist}
% span operator
\DeclareMathOperator{\spn}{span}
% error function
\DeclareMathOperator{\erf}{erf}

% quantities that show up regularly
\newcommand{\wsln}{\ensuremath{\omega(\sqrt{\log n})}}
\newcommand{\wslm}{\ensuremath{\omega(\sqrt{\log m})}}

% support algorithms
\newcommand{\sampleZ}{\algo{Sample}\Z}
\newcommand{\sampleD}{\algo{SampleD}}
\newcommand{\tobasis}{\algo{ToBasis}}
\newcommand{\invert}{\algo{Invert}}
\newcommand{\extbasis}{\algo{ExtBasis}}
\newcommand{\extlattice}{\algo{ExtLattice}}
\newcommand{\randbasis}{\algo{RandBasis}}
\newcommand{\noisy}{\algo{Noisy}}

% problems related to lattices
\newcommand{\svp}{\problem{SVP}}
\newcommand{\gapsvp}{\problem{GapSVP}}
\newcommand{\cogapsvp}{\problem{coGapSVP}}
\newcommand{\usvp}{\problem{uSVP}}
\newcommand{\sivp}{\problem{SIVP}}
\newcommand{\gapsivp}{\problem{GapSIVP}}
\newcommand{\cvp}{\problem{CVP}}
\newcommand{\gapcvp}{\problem{GapCVP}}
\newcommand{\cvpp}{\problem{CVPP}}
\newcommand{\gapcvpp}{\problem{GapCVPP}}
\newcommand{\bdd}{\problem{BDD}}
\newcommand{\gdd}{\problem{GDD}}
\newcommand{\add}{\problem{ADD}}
\newcommand{\incgdd}{\problem{IncGDD}}
\newcommand{\incivd}{\problem{IncIVD}}
\newcommand{\crp}{\problem{CRP}}
\newcommand{\gapcrp}{\problem{GapCRP}}
\newcommand{\dgs}{\problem{DGS}}
% avg-case stuff
\newcommand{\sis}{\problem{SIS}}
\newcommand{\isis}{\problem{ISIS}}
\newcommand{\ilwe}{\problem{ILWE}}
\newcommand{\lwe}{\problem{LWE}}
\newcommand{\rlwe}{\problem{RLWE}}
\newcommand{\declwe}{\problem{WDLWE}}
\newcommand{\extlwe}{\problem{extLWE}}
\newcommand{\dlwe}{\problem{DLWE}}
\newcommand{\lpn}{\problem{LPN}}
\newcommand{\Psibar}{\ensuremath{\bar{\Psi}}}

% set of hints for ExtLWE
\newcommand{\calZ}{\mathcal{Z}}

%% ALGEBRAIC NUMBER THEORY NOTATION

\usepackage{amsbsy}
\usepackage[mathscr]{eucal}

% number field
\newcommand{\nf}{K}
% family of number fields
\newcommand{\nffam}{\mathcal{K}}
% algebraic integers of a number field
\newcommand{\nfints}{\mathcal{O}_{\nf}}
% the norm of an element or ideal
\DeclareMathOperator{\norm}{N}
% trace
\DeclareMathOperator{\trace}{Tr}
% discriminant of number field
\newcommand{\disc}{\Delta}
% different ideal
\newcommand{\diff}{\delta}
% the range of the canonical embedding
\newcommand{\canonrange}{\ensuremath{\R^{s_1} \times \C^{2s_2}}}
% root discriminant
\newcommand{\rootdisc}{\mathscr{D}}
% an ideal given by generators
\DeclarePairedDelimiter\ideal{\langle}{\rangle}
% a frak'ed ideal
\newcommand{\frideal}[1]{\mathfrak{#1}}
% dummy ideals
\newcommand{\I}{\mathcal{I}}
\newcommand{\J}{\mathcal{J}}
\newcommand{\M}{\mathcal{M}}
% radical
\DeclareMathOperator{\rad}{rad}

% crypto-related notation

% KEYS AND RELATED

\newcommand{\mm}[1]{\ensuremath{#1}}

\newcommand{\pp}{\mm{pp}}      % public params
\newcommand{\pk}{\mm{pk}}
\newcommand{\vk}{\mm{vk}}
\newcommand{\sk}{\mm{sk}}
\newcommand{\mpk}{\mm{mpk}}
\newcommand{\msk}{\mm{msk}}
\newcommand{\fk}{\mm{fk}}
\newcommand{\id}{id}
\newcommand{\msgspace}{\ensuremath{\mathcal{M}}}
\newcommand{\idspace}{\ensuremath{\mathcal{ID}}}

\newcommand{\concat}{\ensuremath{\|}}

% GAMES STUFF

% advantage
\newcommand{\advan}{\ensuremath{\mathbf{Adv}}}

% different attack models
\newcommand{\attack}[1]{\ensuremath{\text{#1}}}

\newcommand{\atk}{\attack{atk}} % dummy attack
\newcommand{\indcpa}{\attack{ind-cpa}}
\newcommand{\indcca}{\attack{ind-cca}}
\newcommand{\anocpa}{\attack{ano-cpa}} % anonymous
\newcommand{\anocca}{\attack{ano-cca}}
\newcommand{\sidindcpa}{\attack{sid-ind-cpa}} % selective-id
\newcommand{\aidindcpa}{\attack{aid-ind-cpa}} % adaptive-id
\newcommand{\sidindcca}{\attack{sid-ind-cca}}
\newcommand{\aidindcca}{\attack{aid-ind-cca}}
\newcommand{\euacma}{\attack{eu-acma}} % forgery: adaptive chosen-message
\newcommand{\euscma}{\attack{eu-scma}} % forgery: static chosen-message
\newcommand{\suacma}{\attack{su-acma}} % strong forgery: adaptive chosen-message
\newcommand{\suscma}{\attack{su-scma}} % strong forgery: static chosen-message

% ADVERSARIES
\newcommand{\attacker}[1]{\ensuremath{\mathcal{#1}}}

\newcommand{\Adv}{\attacker{A}}
\newcommand{\Sim}{\attacker{S}}
\newcommand{\Ora}{\attacker{O}}
\newcommand{\Inv}{\attacker{I}}
\newcommand{\For}{\attacker{F}}

% CRYPTO SCHEMES

\newcommand{\scheme}[1]{\ensuremath{\text{#1}}}

% secret-key cryptosystem
\newcommand{\skc}{\scheme{SKC}}
\newcommand{\skcgen}{\algo{Gen}}
\newcommand{\skcenc}{\algo{Enc}} % can also use \kemenc and \kemdec
\newcommand{\skcdec}{\algo{Dec}}

% public-key cryptosystem
\newcommand{\pkc}{\scheme{PKC}}
\newcommand{\pkcsetup}{\algo{Setup}}
\newcommand{\pkcgen}{\algo{Gen}}
\newcommand{\pkcenc}{\algo{Enc}} % can also use \kemenc and \kemdec
\newcommand{\pkcdec}{\algo{Dec}}
\newcommand{\msgsp}{\mathcal{M}}

% extra deniable algorithms for pkc
\newcommand{\biden}{\scheme{BI\mbox{-}DEN}}
\newcommand{\den}{\scheme{DEN}}
\newcommand{\dengen}{\algo{DenGen}} 
\newcommand{\denenc}{\algo{DenEnc}} % can also use \kemenc and \kemdec
\newcommand{\rfake}{\algo{RecFake}}
\newcommand{\sfake}{\algo{SendFake}}

% plan ahead deniability 
\newcommand{\pden}{\scheme{PL\mbox{-}DEN}}
\newcommand{\pdenenc}{\algo{PADenEnc}} % can also use \kemenc and \kemdec
\newcommand{\prfake}{\algo{PARecFake}}
\newcommand{\psfake}{\algo{PASendFake}}

% bitranslucent sets
\newcommand{\bts}{\scheme{BTS}}
\newcommand{\btsgen}{\algo{Gen}}
\newcommand{\btsdengen}{\algo{DenGen}}
\newcommand{\sampleP}{\algo{SampleP}} 
\newcommand{\sampleU}{\algo{SampleU}} 
\newcommand{\testP}{\algo{TestP}}
\newcommand{\fakeScoins}{\algo{FakeSCoins}}
\newcommand{\fakeRcoins}{\algo{FakeRCoins}}

% digital signatures
\newcommand{\sig}{\scheme{SIG}}
\newcommand{\siggen}{\algo{Gen}}
\newcommand{\sigsign}{\algo{Sign}}
\newcommand{\sigver}{\algo{Ver}}

% key-encapsulation mechanism
\newcommand{\kem}{\scheme{KEM}}
\newcommand{\kemgen}{\algo{Gen}}
\newcommand{\kemenc}{\algo{Encaps}}
\newcommand{\kemdec}{\algo{Decaps}}

% symmetric cipher
\newcommand{\sym}{\scheme{SYM}}
\newcommand{\symenc}{\algo{E}}
\newcommand{\symdec}{\algo{D}}

% identity-based encryption
\newcommand{\ibe}{\scheme{IBE}}
\newcommand{\ibesetup}{\algo{Setup}}
\newcommand{\ibeext}{\algo{Ext}}
\newcommand{\ibeenc}{\algo{Enc}}
\newcommand{\ibedec}{\algo{Dec}}

% hierarchical IBE (as key encapsulation)
\newcommand{\hibe}{\scheme{HIBE}}
\newcommand{\hibesetup}{\algo{Setup}}
\newcommand{\hibeext}{\algo{Extract}}
\newcommand{\hibeenc}{\algo{Encaps}}
\newcommand{\hibedec}{\algo{Decaps}}

% binary tree encryption (as key encapsulation)
\newcommand{\bte}{\scheme{BTE}}
\newcommand{\btesetup}{\algo{Setup}}
\newcommand{\bteext}{\algo{Extract}}
\newcommand{\bteenc}{\algo{Encaps}}
\newcommand{\btedec}{\algo{Decaps}}

% trapdoor functions
\newcommand{\tdf}{\scheme{TDF}}
\newcommand{\tdfgen}{\algo{Gen}}
\newcommand{\tdfeval}{\algo{Eval}}
\newcommand{\tdfinv}{\algo{Invert}}
\newcommand{\tdfver}{\algo{Ver}}

% preimage samplable functions
\newcommand{\psf}{\scheme{PSF}}
\newcommand{\psfgen}{\algo{Gen}}
\newcommand{\psfsamdom}{\algo{SampleDom}}
\newcommand{\psfsampre}{\algo{SamplePre}}

% COMPLEXITY MEASURES

\newcommand{\qhash}{\ensuremath{Q_{H}}}
\newcommand{\qext}{\ensuremath{Q_{E}}}
\newcommand{\qid}{\ensuremath{Q_{\text{id}}}}
\newcommand{\qsign}{\ensuremath{Q_{\text{sign}}}}

% UC STUFF

\newcommand{\func}[1]{\ensuremath{\mathcal{F}_{#1}}\xspace}
\newcommand{\prot}[1]{\ensuremath{\mathcal{\pi}_{#1}}\xspace}
\newcommand{\exec}{\ensuremath{\text{EXEC}}}
%\newcommand{\ideal}{\ensuremath{\text{IDEAL}}}
% \Sim and \Adv are already defined above
\newcommand{\Env}{\attacker{Z}}
\newcommand{\sid}{\ensuremath{\mathsf{sid}}\xspace}
\newcommand{\ssid}{\ensuremath{\mathsf{ssid}}\xspace}
\newcommand{\command}[1]{\ensuremath{\mathsf{#1}}\xspace}
% MARGIN NOTES

%\newif\ifnotes\notestrue
\newif\ifnotes\notesfalse


\ifnotes
\usepackage{color}
\definecolor{mygrey}{gray}{0.50}
\newcommand{\notename}[2]{{\textcolor{mygrey}{\footnotesize{\bf (#1:} {#2}{\bf ) }}}}
\newcommand{\noteswarning}{{\begin{center} {\Large WARNING: NOTES ON}\end{center}}}

\else

\newcommand{\notename}[2]{{}}
\newcommand{\noteswarning}{{}}

\fi

\newcommand{\cnote}[1]{{\notename{Chris}{#1}}}
\newcommand{\znote}[1]{{\notename{Zvika}{#1}}}
\newcommand{\onote}[1]{{\notename{Oded}{#1}}}
\newcommand{\anote}[1]{{\notename{Adeline}{#1}}}
\newcommand{\dnote}[1]{{\notename{Damien}{#1}}}






\newcommand{\etal}{{et~al.}}

% new epsilons
\newcommand{\eps}{\varepsilon}
\newcommand{\bfeps}{\boldsymbol{\varepsilon}}
\renewcommand{\epsilon}{\varepsilon}

% paragraph macro (useful because many classes redefine it}
\newcommand{\myparagraph}[1]{\paragraph{{#1}.}}

% Expectations
\DeclareMathOperator{\Exp}{{\mathbb{E}}}

% Macros from Adeline and Damien 
\newcommand{\mult}{\cdot}
\newcommand{\wn}{\omega( \sqrt{\log n})}

% macros for extLWE proof
\newcommand{\roundd}{{\mathrm{round}}}
\newcommand{\fract}{{\mathrm{frac}}}


% Macros for decision problems and such
\def\binset{\{0,1\}}
\def\bbZ{{\mathbb Z}}
\def\bbN{{\mathbb N}}
\def\bbT{{\mathbb T}}
\def\cA{{\cal A}}
\def\cB{{\cal B}}
\def\cZ{{\cal Z}}

%\def\getsr{\stackrel{\scriptscriptstyle{\$}}{\gets}}
\def\getsr{\gets}
\def\getsd{{:=}}

\newcommand{\mx}[1]{\mathbf{{#1}}}
\newcommand{\vc}[1]{\mathbf{{#1}}}
%\newcommand{\gvc}[1]{\bm{{#1}}}
\newcommand{\gvc}[1]{\vc{{#1}}}

\newcommand{\adv}{\mathrm{Adv}}

\newcommand{\zolwe}{\mathsf{binLWE}}
\def\mAh{\hat{\mx{A}}}


% OR: a hack to disable warnings
\renewcommand{\notename}[2]{{}}
\renewcommand{\noteswarning}{{}}

\renewcommand{\vec}{\mathbf}
\newcommand{\calD}{\mathcal{D}}

%%%%

% texlate: fit complete oversized float boxes v1
\usepackage{graphicx}
\begingroup
\makeatletter
\AtBeginDocument{%
\let\texlate@endfloatbox\@endfloatbox
\def\@endfloatbox{%
\texlate@endfloatbox
\def\texlate@figure{figure}%
\def\texlate@figurestar{figure*}%
\def\texlate@table{table}%
\def\texlate@tablestar{table*}%
\let\texlate@floatscope\@empty
\ifx\@currenvir\texlate@figure\def\texlate@floatscope{1}\fi
\ifx\@currenvir\texlate@figurestar\def\texlate@floatscope{1}\fi
\ifx\@currenvir\texlate@table\def\texlate@floatscope{1}\fi
\ifx\@currenvir\texlate@tablestar\def\texlate@floatscope{1}\fi
\ifx\texlate@floatscope\@empty\else
\ifdim\dimexpr\ht\@currbox+\dp\@currbox\relax>\textheight
\edef\texlate@floatwidth{\the\wd\@currbox}%
\edef\texlate@floatheight{\the\dimexpr\textheight-\baselineskip\relax}%
\ifdim\texlate@floatheight>0pt
\typeout{TeXlate-Float-Fit: \@captype\space \csname the\@captype\endcsname; height \the\dimexpr\ht\@currbox+\dp\@currbox\relax; limit \texlate@floatheight}%
\global\setbox\@currbox=\vbox{\hbox to\texlate@floatwidth{\hfil\resizebox*{!}{\texlate@floatheight}{\box\@currbox}\hfil}}%
\fi\fi\fi}%
}
\endgroup
\begin{document}

\title{这是译文}
\subtitle{[这是译文]}
%\titlenote{A full version of this paper is available as
%\textit{Author's Guide to Preparing ACM SIG Proceedings Using
%\LaTeX$2_\epsilon$\ and BibTeX} at
%\texttt{www.acm.org/eaddress.htm}}}



\numberofauthors{5}

\author{
\alignauthor Zvika Brakerski\titlenote{Supported by a Simons Postdoctoral Fellowship and DARPA.}\\
	\affaddr{Stanford University} \\
%	\email{zvika@stanford.edu }
%
\alignauthor Adeline Langlois \\
	\affaddr{Laboratoire LIP, ENS de Lyon
	%\titlenote{(U.\ Lyon, CNRS, ENS Lyon, INRIA, UCBL),
    %46 All\'ee d'Italie, 69364 Lyon Cedex 07, France.}
    }\\
%    \email{adeline.langlois@ens-lyon.fr}
%
\alignauthor Chris Peikert\titlenote{  This material is based upon work supported by the National Science
  Foundation under CAREER Award~CCF-1054495, by DARPA under
  agreement number FA8750-11-C-0096, and by the Alfred P.~Sloan
  Foundation.  Any opinions, findings, and conclusions or
  recommendations expressed in this material are those of the
  author(s) and do not necessarily reflect the views of the National
  Science Foundation, DARPA or the U.S.~Government, or the Sloan
  Foundation.  The U.S. Government is authorized to reproduce and
  distribute reprints for Governmental purposes notwithstanding any
  copyright notation thereon.} \\
	\affaddr{School of Computer Science, Georgia Institute of Technology}
    %  \email{cpeikert@cc.gatech.edu}
%
  \and
%
\alignauthor Oded Regev\titlenote{Supported by a European Research Council (ERC)
  Starting Grant. Part of the work done while the author was with
  the CNRS, DI, ENS, Paris.}\\
  	\affaddr{Courant Institute,}\\
 	\affaddr{New York University}
  %
\alignauthor Damien Stehl\'e \titlenote{The author was
  partly supported by the Australian Research Council Discovery Grant
  DP110100628.} \\
	\affaddr{Laboratoire LIP, ENS de Lyon
	% (U.\ Lyon, CNRS, ENS Lyon, INRIA, UCBL), 46 All\'ee d'Italie, 69364 Lyon Cedex 07, France
 	} \\
%    \email{damien.stehle@ens-lyon.fr}
}




\maketitle

\begin{abstract}
  这是译文 (这是译文) 这是译文 \emph{这是译文} 这是译文 \emph{这是译文} 这是译文.

这是译文.




%%% Local Variables:
%%% mode: latex
%%% TeX-master: "lwehard"
%%% End:

\end{abstract}

%%%% DON'T REMOVE %%%%%
\noteswarning
%%%% DON'T REMOVE %%%%%


\section{这是译文}
\label{sec:introduction}




这是译文 \emph{这是译文} 这是译文 \emph{这是译文} 这是译文.

这是译文 \emph{这是译文} ($\sis$) 这是译文 \emph{这是译文} ($\lwe$) 这是译文~\cite{Ajtai96}, 这是译文 (这是译文) 这是译文 $\gapsvp$, 这是译文 (这是译文.,~\cite{DBLP:journals/siamcomp/MicciancioR07}), 这是译文 $\sis$ 这是译文 $\sis$ 这是译文~\cite{Ajtai96} 这是译文~\cite{DBLP:journals/eccc/ECCC-TR96-042}, 这是译文~\cite{DBLP:conf/crypto/MicciancioV03,DBLP:conf/pkc/Lyubashevsky08,DBLP:conf/asiacrypt/KawachiTX08}, 这是译文~\cite{DBLP:conf/stoc/GentryPV08,DBLP:conf/eurocrypt/CashHKP10,DBLP:conf/pkc/Boyen10,DBLP:conf/eurocrypt/MicciancioP12,DBLP:conf/eurocrypt/Lyubashevsky12}. 

这是译文: \[
\parens{\parens{ \veca_i, \inner{\vc{a}_i, \vc{s}}+e_i \bmod q}}_{i}
\quad \text{and} \quad
\parens{\parens{ \veca_i, u_i }}_{i}~,
\] 这是译文 $\vecs$ 这是译文 $\Z_q^n$ 这是译文 $\veca_i \in \Z_q^n$,  $u_i$ 这是译文 $\Z_q$, 这是译文'' $e_i \in \Z$ 这是译文 (这是译文) 这是译文 $\alpha q$ 这是译文 $\alpha = o(1)$.
%\dnote{I replaced $\ll$ by $o(\cdot)$}


这是译文 $\lwe$ 这是译文~\cite{DBLP:journals/jacm/Regev09,DBLP:conf/crypto/PeikertVW08,DBLP:conf/ctrsa/LindnerP11} 这是译文~\cite{DBLP:conf/stoc/PeikertW08,DBLP:conf/stoc/Peikert09,DBLP:conf/eurocrypt/MicciancioP12} 这是译文~\cite{DBLP:conf/crypto/PeikertVW08}, 这是译文~\cite{DBLP:conf/stoc/GentryPV08,DBLP:conf/eurocrypt/CashHKP10,DBLP:conf/eurocrypt/AgrawalBB10,DBLP:conf/crypto/AgrawalBB10}, 这是译文 (这是译文., \cite{DBLP:conf/tcc/AkaviaGV09,DBLP:conf/crypto/ApplebaumCPS09,DBLP:conf/innovations/GoldwasserKPV10}), 这是译文~\cite{DBLP:conf/focs/BrakerskiV11,DBLP:conf/innovations/BrakerskiGV12,DBLP:conf/crypto/Brakerski12} (这是译文~\cite{DBLP:conf/stoc/Gentry09}), 这是译文~\cite{DBLP:conf/focs/KlivansS06}.

这是译文 $\sis$ 这是译文 $\lwe$ 这是译文 $\lwe$~\cite{DBLP:journals/jacm/Regev09} 这是译文 $\sis$ 这是译文 \emph{这是译文}. 这是译文 $\lwe$, 这是译文 (这是译文) 这是译文 \emph{这是译文} 这是译文 \emph{这是译文} 这是译文 $\lwe$.

这是译文~\cite{DBLP:conf/stoc/Peikert09} (这是译文~\cite{DBLP:conf/crypto/LyubashevskyM09}).  这是译文 $\lwe$ 这是译文 \emph{这是译文} 这是译文 $\lwe$ 这是译文~\cite{DBLP:conf/focs/KlivansS06}, 这是译文 (这是译文) 这是译文 $\lwe$ 这是译文 $2$ (这是译文), 这是译文~\cite{DBLP:conf/stoc/Peikert09} 这是译文 $\lwe$ 这是译文 $\lwe$ 这是译文.


这是译文 $\lwe$ 这是译文 \emph{这是译文} 这是译文 $\lwe$ 这是译文.

\myparagraph{Main result} 这是译文 $\lwe$ 这是译文 $\lwe$ 这是译文.

\onote{should we add somewhere a formal statement of our main result? at least a combination of sections 3 and 4?}
\begin{theorem}[这是译文]\label{thm:informal} 这是译文 $n$-这是译文 $\lwe$ 这是译文 $\poly(n)$ 这是译文 $\sqrt{n}$.
\end{theorem}
这是译文 $\lwe$ 这是译文.

\myparagraph{Techniques} 这是译文.

这是译文 $\lwe$ 这是译文 $q$ 这是译文 $n$ 这是译文 $\lwe$ 这是译文 (这是译文) 这是译文 $p=\poly(n)$ 这是译文 $n \log_2 q$. 这是译文~\ref{thm:informal} 这是译文~\cite{DBLP:conf/stoc/Peikert09}, 这是译文 $q=2^n$ 这是译文 $n$-这是译文 $\gapsvp$. 这是译文 $n$ 这是译文 $n \log_2 q$ 这是译文 (这是译文).%invariant

这是译文 $\Z_q$ 这是译文 $\Z_p$ 这是译文 $a \in \{0,\ldots,q-1\}$ 这是译文 $\floor{pa/q} \in \{0,\ldots,p-1\}$. 这是译文 $\lwe$, 这是译文 (这是译文~\ref{sec:modexpansion}). 这是译文 $\lwe$ 这是译文 $\vecs$, 这是译文 $\vecs$ 这是译文 (这是译文~\ref{sec:shortsecret}) 这是译文 $\lwe$ (这是译文 $n$) 这是译文 $\Z_q^n$ 这是译文 $\lwe$ (这是译文 $n \log_2 q$) 这是译文 $\{0,1\}$. 这是译文 (这是译文~\ref{sec:theshortsecretreduction}) 这是译文-$\lwe$'' 这是译文 $\lwe$ 这是译文 (这是译文~\ref{sec:extlwe}).


这是译文~\cite{DBLP:conf/focs/BrakerskiV11,DBLP:conf/innovations/BrakerskiGV12,DBLP:conf/crypto/Brakerski12}. 这是译文 $\lwe$ 这是译文.


%
这是译文 $\lwe$ 这是译文 $\{0,1\}$ 这是译文.~\cite{DBLP:conf/innovations/GoldwasserKPV10} 这是译文 $\lwe$ 这是译文 $\{0,1\}$-$\lwe$ 这是译文~\cite{DBLP:conf/innovations/GoldwasserKPV10} (这是译文), 这是译文 $\{0,1\}$-$\lwe$ 这是译文, $\lwe$ 这是译文 $\{0,1\}$-$\lwe$ 这是译文.


这是译文~\cite{DBLP:conf/innovations/GoldwasserKPV10} 这是译文 $\lwe$ 这是译文 $\vecz$ 这是译文~\cite{DBLP:conf/crypto/ONeillPW11} 这是译文 (这是译文)这是译文~\cite{DBLP:conf/pkc/Alperin-SheriffP12} 这是译文.



\myparagraph{Broader perspective} 这是译文~\ref{thm:informal}, 这是译文 $\lwe$. 这是译文~\ref{sec:modexpansion} 这是译文 $\lwe$ 这是译文 $n$ 这是译文 $p$ 这是译文 $\lwe$ 这是译文 $n/k$ 这是译文 $p^k$ (这是译文~\ref{cor:mod-dim-tradeoff}). 这是译文 $n$-这是译文 $\lwe$ 这是译文 $q$ 这是译文 $n \log_2 q$. 这是译文 $n$ 这是译文 $q$ 这是译文 $n \log_2 q$ 这是译文 $\lwe$.

这是译文 (这是译文 $n \log_2 q$ 这是译文), 这是译文 $n$-这是译文 $\lwe$ 这是译文 $2^n$ 这是译文 $n^2$-这是译文 $\lwe$ 这是译文, $n$-这是译文 $\lwe$ 这是译文 $2^n$, 这是译文~\cite{DBLP:conf/stoc/Peikert09} 这是译文 $n$-这是译文 $n^2$-这是译文 $\exp(\widetilde\Omega(n^2))$ 这是译文~\cite{DBLP:conf/stoc/Peikert09} (这是译文~\ref{thm:informal}) 这是译文~\cite{DBLP:journals/jacm/Regev09}.

\onote{for the final version, consider adding precise statements for this and the next point} 这是译文 $1$-这是译文 $\lwe$ 这是译文 $2^n$ 这是译文 $n$-这是译文 $\lwe$ 这是译文 $1$-这是译文 $\lwe$ 这是译文~\cite{DBLP:conf/crypto/BonehV96}. 这是译文~\cite{DBLP:conf/stoc/AjtaiD97} 这是译文~\cite{DBLP:journals/jacm/Regev04}, 这是译文~\cite{RegevSurvey}. 这是译文 $1$-这是译文-$\lwe$ 这是译文~\cite{DBLP:conf/eurocrypt/LyubashevskyPR10} (这是译文~1, 这是译文~$\Z$). 这是译文~\cite{DBLP:conf/scn/GentryHPS12}, 这是译文-$\lwe$, 这是译文 (这是译文). 这是译文-$\sis$ 这是译文~\cite{DBLP:conf/icalp/LyubashevskyM06,DBLP:conf/tcc/PeikertR06} 这是译文 $\lwe$-这是译文-$\sis$ 这是译文. (这是译文 $1$, 这是译文.)

这是译文 (这是译文) \lwe 这是译文~\ref{cor:mod-reduction-2}.  这是译文 (这是译文.,~\cite{DBLP:journals/jacm/Regev09,DBLP:conf/stoc/Peikert09,DBLP:conf/crypto/MicciancioM11,DBLP:conf/eurocrypt/MicciancioP12}) 这是译文 \emph{这是译文}, 这是译文.

\onote{for final version: decide what to do with real-LWE}

\onote{for final version, consider applying these ideas to SIS, especially the 1-dim case which is like subset sum, and more generally, modulus/dimension tradeoff for SIS. Chris said:
The idea is to do a similar transformation as in Lemma 4.4, but on SIS samples, going from a in $\Z_q^n$ to $a' \in \Z_{q'}^{n'}$.  Then given a good SIS solution $x$ to $A'x = 0$, we also have $G^T A' x = 0$ and therefore $A x \approx 0$.  Thus $x$ is a solution to the original SIS instance, because it suffices to find a preimage of "nearly" 0.  This is because if $Ax = e \approx 0$, then we get the homogeneous solution $[I | A] (e; x) = 0$, and the instance $[I | A]$ is just SIS in (hermite) normal form.}

\dnote{Something that I find interesting is that the complexity of the
  best known attack is~$2^{O(n \log q/ \log^2 \alpha)}$, which is
  constant wrt $n \log q$. This was not reflected by our previous
  understanding of the complexity of \lwe, and we manage to close the
  gap between the hardness results and the algorithms. By the way,
  this~$2^{O(n \log q/ \log^2 \alpha)}$ bound suggests it may be
  possible to transfer some of~$n$ and~$\log q$ into~$\log \alpha$. Do
  we know \lwe self-reductions from, say, $q,\alpha$ to~$p,\beta$ such
  that~$\log q/ |\log \alpha| = \log p/|\log \beta|$? or from
  $n,\alpha$ to~$n',\beta$ such that~$n/ |\log \alpha| = n'/|\log
  \beta|$}

\onote{maybe mention Ajtai's 2005 paper using Dirichlet's problem?}

\myparagraph{Open questions} 这是译文~\ref{thm:informal} 这是译文~\cite{DBLP:conf/stoc/Peikert09} 这是译文~\cite{DBLP:journals/jacm/Regev09} 这是译文 $\sis$.  这是译文~\cite{DBLP:conf/stoc/Peikert09} 这是译文~\cite{DBLP:journals/jacm/Regev09} 这是译文 $\sis$ 这是译文~\ref{thm:informal}  这是译文~\cite{DBLP:conf/stoc/Peikert09} 这是译文 \gapsvp (这是译文 \bdd 这是译文 \usvp~\cite{DBLP:conf/crypto/LyubashevskyM09}), 这是译文 \sivp, 这是译文~\cite{DBLP:journals/jacm/Regev09} 这是译文 \sis 这是译文-\lwe, 这是译文 \gapsvp 这是译文. 

这是译文~\cite{DBLP:conf/stoc/Peikert09} 这是译文~\ref{thm:informal}. 这是译文 $n$-这是译文 $\lwe$ 这是译文 $n$-这是译文 (这是译文 $\sqrt{n}$-这是译文)? 这是译文~\cite{DBLP:journals/jacm/Regev09}.

这是译文 (这是译文) 这是译文 $\lwe$ 这是译文 $2^{O(\sqrt{n})}$. 这是译文~\ref{thm:informal}, 这是译文 $2^{O(\sqrt{n})}$-这是译文 \emph{这是译文} 这是译文~\cite{DBLP:journals/jacm/Regev09}. 这是译文 $2^{O(\sqrt{n})}$-这是译文~\cite{DBLP:journals/siamcomp/Kuperberg05} (这是译文~\cite{DBLP:journals/siamcomp/Regev04}).




%%% Local Variables:
%%% mode: latex
%%% TeX-master: "lwehard"
%%% End:


\section{这是译文}
\label{sec:preliminaries}

这是译文 $\T =
\R/\Z$ 这是译文 $1$. 这是译文 $\T_q$ 这是译文 $q$, 这是译文 $\{0,1/q,\ldots,(q-1)/q\}$.

这是译文 $P,Q$ 这是译文 $\sum |P(i)-Q(i)|/2$ 这是译文~$i$ 这是译文 (这是译文.,~\cite[Eq.~(2.3)]{AldousDiaconis87} 这是译文).
\begin{claim}\label{clm:separationdist}  这是译文 $P$ 这是译文 $Q$ 这是译文~$P(i) \ge
  (1-\eps)Q(i)$ 这是译文~$i$, 这是译文 $P$ 这是译文 $Q$ 这是译文~$\eps$.
\end{claim}



这是译文~\cite{DBLP:journals/siamcomp/HastadILL99}.

\begin{lemma}
  \label{lem:lhl}  这是译文 $k, n, q \ge 1$ 这是译文 $\epsilon > 0$ 这是译文 $n
  \ge k \log_2 q + 2 \log_2 (1/\epsilon)$. 这是译文 $\mx{H} \gets \T_q^{k \times
    n}$, $\vc{z} \gets \binset^n$, $\vc{u} \gets \bbT_q^k$, 这是译文 $(\mx{H}, \mx{H}\vc{z})$ 这是译文 $(\mx{H}, \vc{u})$ 这是译文 $\epsilon$.
\end{lemma}


这是译文 \emph{这是译文} $P$ 这是译文 $P_0$ 这是译文 $P_1$, 这是译文 \emph{这是译文} 这是译文 $\cA$ 这是译文 $P$ 这是译文 \[
\adv[\cA] = \abs{\Pr[\cA(P_0)] - \Pr[\cA(P_1)]}~.
\] 这是译文~$P$ 这是译文~$Q$ 这是译文 (这是译文) 这是译文~$\cA^{\cB}$ 这是译文~$P$ 这是译文~$\cB$ 这是译文~$Q$. 这是译文 (这是译文~\ref{lem:lelwe}) 这是译文. 


\subsection{这是译文}
\label{sec:lattices}

这是译文 $n$-这是译文 (这是译文) 这是译文 $\Lambda \subseteq \R^n$ 这是译文 $n$ 这是译文 \emph{这是译文} 这是译文 $\matB = \set{\vecb_{1}, \ldots,
  \vecb_{n}} \subseteq \R^n$, \[ \Lambda = \lat(\matB) = \set[\Big]{
  \sum_{i \in [n]} z_i \vecb_i\; :\; \vecz \in \mathbb{Z}^n }. \] 这是译文 \emph{这是译文} 这是译文 $\Lambda \subset \R^n$ 这是译文 $\Lambda^* = \set{ \vecx \in \R^n : \inner{\Lambda,\vecx} \subseteq \Z}$.

这是译文 \emph{这是译文} (这是译文 \emph{这是译文}) $\lambda_1(\Lambda)$ 这是译文 $\Lambda$ 这是译文., $\lambda_1(\Lambda) = \min_{\veczero \neq
  \vecx \in \Lambda} \length{\vecx}$.  这是译文 $\gamma = \gamma(n) \ge 1$, 这是译文 $\gapsvp_\gamma$ 这是译文 $\matB$ 这是译文 $n$-这是译文 $\Lambda = \lat(\matB)$ 这是译文 $d$, 这是译文 $\lambda_1(\lat(\matB)) \le d$  这是译文 $\lambda_1(\lat(\matB)) > \gamma d$. 这是译文~\cite{KhotLLL25,RegevLLL25} 这是译文 $\gapsvp_\gamma$. \dnote{I added that sentence}


\subsection{这是译文}
\label{sec:gaussian_measures}


这是译文 $r > 0$, 这是译文 $n$-这是译文 $\rho_r : \R^{n} \to (0,1]$ 这是译文 \[ \rho_r(\vecx) := \exp(-\pi
\length{\vecx}^{2}/r^2).
\] 这是译文., $\rho_r(A) = \sum_{\vecx \in A} \rho_r(\vecx) 
\in [0,+\infty]$ 这是译文~$A \subseteq \R^{n}$. 这是译文 (这是译文) 这是译文~$D_r$ 这是译文 $\rho_r$. 这是译文 $\matB$, 这是译文 $D_\matB$ 这是译文 $\matB \vecx$ 这是译文 $\vecx$ 这是译文 $D_1$. 这是译文 $\matB$ 这是译文 \[
\exp(-\pi \vecx^T (\matB \matB^T)^{-1} \vecx).
\] 这是译文 $\matB_1,\matB_2$, 这是译文 $D_{\matB_1}$ 这是译文~$D_{\matB_2}$ 这是译文 $D_{\matC}$ 这是译文 $\matC = (\matB_1^{} \matB_1^T + \matB_2^{} \matB_2^T)^{1/2}$.

这是译文 $n$-这是译文 $\Lambda$ 这是译文 $\vecu \in \R^n$, 这是译文 \emph{这是译文}~$D_{\Lambda+\vecu,r}$ 这是译文 $\Lambda+\vecu$ 这是译文 $\rho_r$. 这是译文 (这是译文) 这是译文  (\cite[Theorem~4.1]{DBLP:conf/stoc/GentryPV08}; 这是译文~\cite{DBLP:conf/crypto/Peikert10}). 这是译文~\ref{sec:exactgpv}, 这是译文 \cite{DBLP:conf/stoc/GentryPV08}. 这是译文 (这是译文) 这是译文.

\begin{lemma}\label{lem:gpv} 这是译文 $\matB$ 这是译文 $n$-这是译文 $\Lambda = \lat(\matB)$, $\vecc \in \R^n$, 这是译文 $r \ge \|\widetilde \matB\| \cdot \sqrt{\ln(2n+4)/\pi}$, 这是译文 $D_{\Lambda+\vecc,r}$. 
\end{lemma}

这是译文, $\widetilde \matB$ 这是译文 $\matB$, 这是译文 $\|\widetilde \matB\|$ 这是译文 \emph{这是译文} 这是译文~\cite{DBLP:journals/siamcomp/MicciancioR07}.

\begin{definition}
  \label{def:smoothing}  这是译文 $\Lambda$ 这是译文 $\epsilon > 0$, 这是译文 $\smootheps(\Lambda)$ 这是译文 $s >
  0$ 这是译文 $\rho_{1/s}(\Lambda^* \setminus \set{\veczero}) \leq
  \epsilon$.
\end{definition}

\begin{lemma}[{{\cite[Lemma~3.1]{DBLP:conf/stoc/GentryPV08}}}]\label{lem:boundonsmoothing} 这是译文 $\eps>0$ 这是译文 $n$-这是译文 $\Lambda$ 这是译文 $\matB$, \[
\smootheps(\Lambda) \le \|\widetilde \matB \| \sqrt{\ln(2n(1+1/\eps))/\pi}.
\]
\end{lemma}


这是译文.

\begin{lemma}[{{\cite[Lemma 4.1]{DBLP:journals/siamcomp/MicciancioR07}}}]
  \label{lem:smoothingcontinuous}  这是译文 $n$-这是译文 $\Lambda$, $\epsilon>0$, $r \geq
  \eta_{\epsilon}(\Lambda)$, 这是译文 $\vecx \bmod \Lambda$ 这是译文 $\vecx \leftarrow D_{r}$ 这是译文 $\eps/2$ 这是译文 $\Lambda$.
\end{lemma}

\begin{lemma}[{{\cite[Claim 3.8]{DBLP:journals/jacm/Regev09}}}]
  \label{lem:smoothing}  这是译文 $n$-这是译文 $\Lambda$, $\epsilon>0$, $r \geq
  \eta_{\epsilon}(\Lambda)$, 这是译文 $\vecc \in \R^n$, 这是译文 $\rho_{r}(\Lambda + \vecc) \in [
  \tfrac{1-\epsilon}{1+\epsilon}, 1 ] \cdot \rho_{r}(\Lambda)$.
\end{lemma}

\begin{lemma}[{{\cite[Claim~3.9]{DBLP:journals/jacm/Regev09}}}]
\label{le:discrete_plus_cont_Gauss} 这是译文~$\Lambda$ 这是译文 $n$-这是译文~$\vec{u}\in \R^n$ 这是译文 $r, s > 0$ 这是译文 $t=\sqrt{r^2+ s^2}$. 这是译文 $rs/t=1/\sqrt{1/r^2+1/s^2} \geq \eta_{\eps} (\Lambda)$ 这是译文~$\eps <
1/2$. 这是译文~$Y$ 这是译文~$\R^n$ 这是译文~$D_{\Lambda+\vec{u},r}$ 这是译文~$D_s$. 这是译文~$Y$ 这是译文~$D_t$ 这是译文~$4\eps$.
\end{lemma}

\begin{lemma}[{{\cite[Corollary~3.10]{DBLP:journals/jacm/Regev09}}}]
\label{lem:smoothip} 这是译文~$\Lambda$ 这是译文 $n$-这是译文~$\vc{u}, \vc{z}\in \R^n$ 这是译文 $r, \alpha > 0$. 这是译文 $\parens{1/r^2+\left(\|{\vc{z}}\|/\alpha\right)^2}^{-1/2} \geq \eta_{\eps} (\Lambda)$ 这是译文~$\eps <
1/2$. 这是译文 $\langle\vc{z},\vc{v}\rangle+e$ 这是译文 $\vc{v} \getsr D_{\Lambda+\vc{u}, r}$ 这是译文 $e \getsr D_{\alpha}$, 这是译文 $4\eps$ 这是译文 $D_{\beta}$ 这是译文 $\beta = \sqrt{(r \|\vc{z}\|)^2+\alpha^2}$.
\end{lemma}

\begin{lemma}[{{Special case of~\cite[Theorem 3.1]{DBLP:conf/crypto/Peikert10}}}]
\label{lem:peikertdiscretegauss} 这是译文 $\Lambda$ 这是译文 $r,s>0$ 这是译文 $s \ge \eta_\eps(\Lambda)$ 这是译文 $\eps \le 1/2$. 这是译文 $\vecx$ 这是译文 $D_{r}$ 这是译文 $\vecy$ 这是译文 $D_{\Lambda-\vecx,s}$ 这是译文 $\vecx+\vecy$ 这是译文 $8 \eps$ 这是译文 $D_{\Lambda,(r^2+s^2)^{1/2}}$.
\end{lemma}


\subsection{这是译文}
\label{sec:computational_problems}

这是译文 $n,q \ge 1$, 这是译文 $\vec{s}
\in \Z^n$, 这是译文 $\phi$ 这是译文~$\R$, 这是译文 $A_{q,\vecs, \phi}$ 这是译文 $\T_q^n \times \T$ 这是译文 $\veca \in \T_q^n$ 这是译文 $e$ 这是译文 $\phi$, 这是译文 $(\veca, b =
\inner{\veca, \vecs}+e) \in \T_q^n \times \T$.

\begin{definition}
  \label{def:lwe}  这是译文 $n,q \ge 1$, 这是译文 $\phi$  这是译文 $\R$, 这是译文 $\calD$ 这是译文 $\Z^n$, 这是译文 (这是译文) 这是译文 $\lwe$  这是译文 $\lwe_{n,q,\phi}(\calD)$, 这是译文 $\T_q^n \times \T$ 这是译文  $A_{q,\vec{s},\phi}$ 这是译文 $\vec{s}$ 这是译文 $\calD$. 这是译文 $m\in\bbN$ 这是译文 $\lwe_{n,m,q,\phi}(\calD)$.
\end{definition}

这是译文 $A_{q,\vec{s},\phi}$ 这是译文 $\vec{s} \bmod q$, 这是译文 $\vec{s} \in \{0,\ldots,q-1\}^n$. 这是译文 $\calD$, 这是译文 $\lwe_{n,q,\phi}(\calD)$ 这是译文 $\lwe_{n,q,\phi}(U(\{0,\ldots,q-1\}^n))$, 这是译文 $\lwe_{n,q,\phi}$ 这是译文 (这是译文).  这是译文 $\alpha>0$, 这是译文., $\phi=D_\alpha$, 这是译文 $\lwe_{n,q,\alpha}(\calD)$. 这是译文 $\calD$ 这是译文 $\{0,1\}^n$ 这是译文 $\zolwe_{n,q,\phi}$ (这是译文~$\zolwe_{n,m,q,\phi}$ 这是译文~$m$ 这是译文). 这是译文 $\lwe$ 这是译文 (这是译文) 这是译文 \lwe, 这是译文~\cite{DBLP:conf/crypto/ApplebaumCPS09} 这是译文 $q$. 这是译文 $q$, 这是译文~\ref{clm:invertiblemodq} 这是译文.

\begin{lemma}\label{lem:normalform}  这是译文 $q \ge 25$, $n,m \geq 1$, $\alpha>0$, $\eps < 1/2$ 这是译文~$s \geq
  \sqrt{\ln(2n(1+1/\epsilon)/\pi)}/q$, 这是译文  (这是译文) 这是译文~$\lwe_{n,m,q,\alpha}$ 这是译文  $\lwe_{n,m',q,\alpha}(\calD)$  这是译文~$m' = m - (16n + 4\ln \ln q)$ 这是译文 $\calD=D_{\Z^n,q(\alpha^2+s^2)^{1/2}}$, 这是译文 $\zeta$ 这是译文 $(\zeta-8\eps)/4$.  这是译文 $\alpha \geq \sqrt{\ln(2n(1+1/\epsilon)/\pi)}/q$, 这是译文 $s=\alpha$,  这是译文 $\calD=D_{\Z^n,\sqrt{2}q\alpha}$.
\end{lemma}

\begin{proof}
  这是译文 $16n + 4\ln\ln q$ 这是译文 $(\vec{a},b)$.  这是译文~\ref{clm:invertiblemodq}, 这是译文  $1-2e^{-1} \ge 1/4$, 这是译文 $A_0 \in \Z_q^{n \times n}$ 这是译文 $\vec{a}$ 这是译文 (这是译文 $q$)  这是译文 $A_0^{-1} \in \Z_q^{n \times n}$ 这是译文 $q$. 这是译文 $\vec{b}_0 \in \T^n$ 这是译文 $b$ 这是译文 $\vec{b}'_0 \in \T_q^n$ 这是译文 $\vec{b}_0 + \vec{x}$ 这是译文 $\vec{x}$ 这是译文 $D_{q^{-1}\Z^n - \vec{b}_0,s}$.  (这是译文 $q^{-1}\Z^n - \vecb_0$ 这是译文 $\vecb_0$ 这是译文 $\Z^{n} \subseteq q^{-1}\Z^n$.)  这是译文 $m'$ 这是译文 $(\vec{a},b) \in \T_q^n \times \T$ 这是译文  \[
	  \big(\vec{a}' = A_0^{-1} \vec{a}, b' = b - \inner{A_0^{-1} \cdot q \vec{a}, \vec{b}'_0}\big) \in \T_q^n \times \T.
	\]  这是译文 $\lwe$ 这是译文 $m'$ 这是译文. 
		
	这是译文 $A_0$ 这是译文 $\vec{a}$ 这是译文 $\lwe$ 这是译文 $A_0$.  这是译文 $(\veca,b)$ 这是译文 $\T_q^n \times \T$  这是译文 $(\veca',b')$.  这是译文 $A_{q,\vec{s},D_{\alpha}}$  这是译文 $\vec{s} \in \Z^n$. 这是译文~$s \geq \eta_{\epsilon}(q^{-1}\Z)$ 这是译文~\ref{lem:boundonsmoothing},  这是译文~\ref{lem:peikertdiscretegauss} 这是译文 $\vec{b}_0' = q^{-1} A_0^T \vec{s} + \vece_0$  这是译文 $\vece_0$ 这是译文 $8\eps$ 这是译文 $D_{q^{-1}\Z^n, (\alpha^2+s^2)^{1/2}}$.  这是译文 $(\veca',b')$ 这是译文  \[
	   b' = b - \inner{A_0^{-1} \cdot q \vec{a}, \vec{b}'_0} 
		    = \inner{\veca,\vecs}+ e - \inner{\veca, \vecs} - \inner{A_0^{-1} q \veca, \vece_0}
				= \inner{-q \vec{e}_0, \veca'}+e,
	\]  这是译文 $e$ 这是译文 $D_\alpha$. 这是译文 $A_{q,-q\vec{e}_0,D_{\alpha}}$, 这是译文.
\end{proof}


\begin{claim}\label{clm:invertiblemodq} 这是译文 $q \ge 25$, $n \ge 1$, 这是译文 $t_1 \ge 4, t_2 \ge 1$, 这是译文 $t_1 n+t_2 \ln\ln q$ 这是译文 $\vec{a}_1,\vec{a}_2,\ldots$ 这是译文 $\Z_q^n$, 这是译文 $e^{-t_1 n /16} + e^{-t_2/4}$, 这是译文 $n$ 这是译文 $n \times n$ 这是译文 $q$.  这是译文. 
\end{claim}
\begin{proof}
这是译文 $k$ 这是译文 $0$, 这是译文 $A \in \Z_q^{n \times k}$ 这是译文~$U \in \Z^{n \times n}$,  这是译文~$U \cdot A \in \Z_q^{n \times k}$  这是译文 $k \times k$ 这是译文~$q$; 这是译文 $(n-k) \times
  k$ 这是译文 $\vec{a}_i$ 这是译文 $\vec{a}$, 这是译文~$n-k$ 这是译文~$U
  \vec{a}$, 这是译文 $g$, 这是译文~$q$, 这是译文 $\vec{a}$ 这是译文 (这是译文) 这是译文 $V$ 这是译文 $k$ 这是译文 $n-k$ 这是译文 $VU\vec{a}$ 这是译文 $(g,0,\ldots,0)$,  这是译文 $U$ 这是译文 $VU$, 这是译文 $k$. 

这是译文 $k=n$. 这是译文~$k<n-1$ 这是译文~$k = n-1$, 这是译文 $q$ 这是译文 $q$ 这是译文  \[
\prod_{p|q,~p~\text{prime}} (1-p^{-2}) \ge \prod_{p~\text{prime}} (1-p^{-2}) = \zeta(2)^{-1} \approx 0.61,
\] 这是译文 $\zeta$ 这是译文 $q$ 这是译文 $q$ 这是译文 $\varphi(q)/q$, 这是译文 $\varphi$ 这是译文~\cite[Theorem 8.8.7]{BachShallit},  这是译文 $(e^\gamma \ln \ln q + 3/(\ln\ln q))^{-1}$ 这是译文 $\gamma$ 这是译文 $q \ge 25$ 这是译文 $(4 \ln \ln q)^{-1}$. 

这是译文 (这是译文) 这是译文 $U\vec{a}$ 这是译文 $\Z_q^n$ 这是译文 $U$ 这是译文 $k < n-1$ 这是译文 $t_1 n$ 这是译文 $e^{-t_1 n /16}$. 这是译文 $k=n-1$, 这是译文 $t_2 \ln \ln q$ 这是译文 $k=n-1$ 这是译文 $e^{-t_2/4}$.
\end{proof}



\myparagraph{Unknown (Bounded) Noise Rate} 这是译文 $\lwe$ 这是译文 $\beta \le \alpha$ (这是译文 $\alpha$), 这是译文 $\beta$ 这是译文 $\vecs$. 这是译文.

\begin{definition}\label{def:lelwe} 这是译文 $n,q \ge 1$ 这是译文 $\alpha \in (0,1)$, $\lwe_{n,q,\le\alpha}$ 这是译文 $\lwe_{n,q,\beta}$ 这是译文 $\beta = \beta(\vc{s}) \le \alpha$.
\end{definition}

\begin{lemma}\label{lem:lelwe} 这是译文 $\cA$ 这是译文 $\lwe_{n, m, q, \alpha}$ 这是译文 $\epsilon>0$. 这是译文 $\cB$ 这是译文 $\lwe_{n,m',q,\le\alpha}$ 这是译文~$\cA$ 这是译文 $1/3$, 这是译文 $m'$ 这是译文 $\poly(m,1/\epsilon, n, \log q)$.
\end{lemma}

这是译文 (这是译文., \cite[Lemma 3.7]{DBLP:journals/jacm/Regev09} 这是译文 $\lwe$). 这是译文 $\cA$'这是译文 $\cA$'这是译文.

\znote{Commented out proof is down here. Fix or remove.
%
%\begin{proof}
%\znote{Only sketch for now. Also I am not sure all the specific constants are correct, need to double check.}
%Let $\cA$ be as in the lemma statement. Note that $\adv[\cA] \ge \epsilon$ means that
%\begin{equation}
%\Pr_{\vc{s}}\left[\abs{\Pr[\cA(\mx{A}, (\mx{A}\cdot\vc{s})/q+\vc{e})] -  \Pr[\cA(\mx{A}, \vc{u})] \ge \epsilon/2 }\right] \ge \epsilon/2~.
%\end{equation}
%Namely, for $\epsilon/2$ fraction of the secret vectors we get a good distinguishing gap.
%
%We define the algorithm $\cB$ as follows: $\cB$ will draw many samples $(\mx{A}, \vc{d})$ from its input distribution (recall that either all are of the form $A_{\vc{s}, \beta}$ or all are uniform). It will add to each of these samples an additional noise from $D_{\gamma}$, for all $\gamma$ over a fine grid of polynomially many values between $0$ and $\alpha$.
%
%Then $\cB$ will estimate $\cA$'s acceptance probability on the uniform distribution and on each of the $\gamma$ distributions. The estimate is such that with probability $1-\epsilon/6$, no estimate deviates by more than $\epsilon/20$ from the correct value. This can be done in $\poly(1/\epsilon, n, T)$ time.
%
%Then, if either of the $\gamma$'s yielded a value that is different from the value of the uniform distribution by more than $\epsilon/8$, then output $1$. Otherwise output $0$.
%
%The analysis proceeds as follows: First, we note that when $\cB$ is given the uniform distribution, it will output $0$ with all but $\epsilon/6$ probability. Next, we analyze $\cB$'s behavior on $A_{\vc{s}, \beta}$.
%
%% consider the distributions $(\mx{A}, \mx{A}^T\vc{s}/q+\vc{e})$ for $\vc{e} \getsr D_{\sqrt{\beta^2+\gamma^2}}$ for a fine grid of $\gamma$ between $0$ and $\alpha$. Note that these distributions are easy to generate without knowing $\beta$.
%
%If we choose the grid of $\gamma$'s fine enough, then there exists $\gamma$ such that $D^m_{(\beta^2+\gamma^2)^{1/2}}$ and $D^m_{\alpha}$ are within $\epsilon/4$ statistical distance. This means that for this specific $\gamma$, the algorithm $\cA$ will have a distinguishing gap $\epsilon/4$ for at least an $\epsilon/2$ fraction of the $\vc{s}$. \znote{Is the statistical distance between two Gaussians bounded by the difference of their standard deviations or the difference of variances?}\onote{use Claim 2.2 of~\cite{DBLP:journals/jacm/Regev09} which bounds the statistical distance between two normal variables by $9(\beta/\alpha - 1)$ assuming $\alpha < \beta \le 2 \alpha$} This means that $\cB$ will output $1$ with probability at least $\epsilon/2 - \epsilon/6 \ge \epsilon/3$.
%
%This implies a distinguishing gap of at least $\epsilon/6$, which can be improved to $\epsilon/3$ if we notice that we double counted the event of erroneous estimates of $\cA$'s acceptance probability.
%\end{proof}
}






\myparagraph{Relation to Lattice Problems} 这是译文~\cite{DBLP:journals/jacm/Regev09} 这是译文~\cite{DBLP:conf/stoc/Peikert09} 这是译文 (这是译文) 这是译文 $\gapsvp$ 这是译文 \onote{added:}这是译文 $\lwe$.  (这是译文~\cite{DBLP:journals/jacm/Regev09} 这是译文 $\sivp$.) 这是译文 $q$ 这是译文~$n$. 这是译文 {{\cite[Theorem~3.1]{DBLP:journals/jacm/Regev09}}} 这是译文 {{\cite[Theorem~3.1]{DBLP:conf/stoc/Peikert09}}}.

\begin{theorem}\label{thm:worstcaseaveragecase} 这是译文 $n, q \ge 1$ 这是译文 $\alpha \in (0,1)$ 这是译文 $\alpha q \ge 2 \sqrt{n}$. 这是译文 $n$-这是译文 $\gapsvp_{\widetilde{O}(n/\alpha)}$ 这是译文 $\lwe_{n,q,\alpha}$. 这是译文 $q \ge 2^{n/2}$ 这是译文. 
\end{theorem}

这是译文 \emph{这是译文} 这是译文 \lwe, 这是译文 (这是译文.,~\cite{DBLP:journals/jacm/Regev09,DBLP:conf/stoc/Peikert09,DBLP:conf/eurocrypt/MicciancioP12}). 这是译文~\cite{DBLP:conf/eurocrypt/MicciancioP12}, 这是译文~$q$ 这是译文~$2$, 这是译文 (这是译文) $\lwe$ 这是译文~$q$, 这是译文~\cite{DBLP:conf/eurocrypt/MicciancioP12}, 这是译文~$q$, 这是译文 (这是译文~\ref{cor:mod-reduction-2}).

\begin{theorem}[{{Special case of~\cite[Theorem 3.1]{DBLP:conf/eurocrypt/MicciancioP12}}}]\label{thm:searchtodecisionmicpei} 这是译文 $q$ 这是译文 $2$, 这是译文 $\alpha$ 这是译文 $1/q < \alpha < 1/\omega(\sqrt{\log n})$. 这是译文 $\lwe_{n,q,\alpha}$ 这是译文 (这是译文) $\lwe_{n,q,\alpha'}$ 这是译文 $\alpha'=\alpha \cdot \omega(\log n)$.
\end{theorem}





%%% Local Variables:
%%% mode: latex
%%% TeX-master: "lwehard"
%%% End:





\section{这是译文}
\label{sec:shortsecret}



这是译文. 

%\dnote{I changed~$\alpha, \delta \in (0,1)$ to~$\alpha,\delta >0$}

\begin{theorem}\label{thm:zolwe} 这是译文~$k, q \ge 1$, 这是译文 $m \ge n \ge 1$ 这是译文 $\epsilon \in (0,1/2)$, $\alpha, \delta >0$, 这是译文 $n \ge (k+1) \log_2 q + 2 \log_2 (1/\delta)$, $\alpha \ge \sqrt{\ln(2n(1+1/\eps))/\pi}/q$. 这是译文 (这是译文) 这是译文 $\lwe_{k, m, q, \alpha}$ 这是译文 $\zolwe_{n, m, q, \le \sqrt{10n} \alpha}$, 这是译文 $\zeta$, 这是译文  \begin{equation}\label{eq:advantagelossinzolwe}
(\zeta-\delta)/(3m) - 41 \eps/2 - \sum_{p|q,~p~\text{prime}} p^{-k-1}\,.
\end{equation}
\end{theorem}

这是译文~\ref{thm:zolwe} 这是译文~\ref{cor:mod-reduction} (这是译文 $\{0,1\}^n$ 这是译文 $(\sqrt{n},0)$ 这是译文), 这是译文 $\zolwe$ 这是译文 $\zolwe_{n,m,q',\beta}$ 这是译文 $q' \ge 1$ 这是译文 $\xi>0$ 这是译文 \[
\beta := \left(10 n \alpha^2 + \frac{4n}{\pi q'^2} \ln(2n(1+1/\xi))\right)^{1/2},
\] 这是译文~\eqref{eq:advantagelossinzolwe} 这是译文 $14 \xi m$. 这是译文  $\lwe$ 这是译文 $k=\sqrt{n}$ 这是译文 $q=2^{k/2}$ (这是译文 $k$ 这是译文) 这是译文 $\sqrt{n}$-这是译文 (这是译文~\ref{thm:worstcaseaveragecase} 这是译文~\ref{thm:searchtodecisionmicpei}),  这是译文~\ref{thm:informal}. 这是译文 $q'$ 这是译文 $\sqrt{n}$.

这是译文 $q$ 这是译文~\eqref{eq:advantagelossinzolwe} 这是译文~\ref{thm:informal} 这是译文 $q$ 这是译文 $2$, 这是译文 $2^{-k-1}=2^{-\sqrt{n}-1}$, 这是译文 (这是译文~\ref{cor:mod-reduction-2}) 这是译文~\ref{thm:zolwe} 这是译文 $q$ 这是译文. (这是译文~\ref{lem:lelwe} 这是译文 $\lwe$ 这是译文~\ref{cor:mod-reduction-2} 这是译文.) 这是译文 $q^{-\sqrt{n}-1}$ 这是译文 $2^{-n}$. 

这是译文~\ref{sec:firstiserrorless}, 这是译文 $\lwe$ 这是译文~\ref{sec:extlwe}, 这是译文 $\extlwe$, 这是译文 $\lwe$ 这是译文~\ref{sec:theshortsecretreduction}, 这是译文 $\extlwe$ 这是译文 $\lwe$ 这是译文 $\{0,1\}$ 这是译文 (这是译文~\ref{fig:summary}). 

\begin{proof}
这是译文 $m \ge n$, 这是译文~\ref{lem:firsterrorlesshard} 这是译文 $\lwe_{k, m, q, \alpha}$ 这是译文 $\lwe_{k+1, n, q, \alpha}$, 这是译文 $2^{-k+1}$.  这是译文~\ref{lem:redtoextlwe} 这是译文 $\cZ = \{0,1\}^n$, 这是译文 $\xi=2$ 这是译文~\ref{clm:qualityofzeroone}, 这是译文 $\extlwe_{k+1, n, q,\sqrt{5}\alpha,\{0,1\}^n}$ 这是译文~$33 \eps/2$. 这是译文~\ref{lem:multiinstanceextlwe} 这是译文 $\extlwe^m_{k+1, n, q, \sqrt{5} \alpha,\{0,1\}^n}$, 这是译文~$m$ 这是译文~\ref{lem:zolwe} 这是译文 $\zolwe_{n, m, q, \le \sqrt{10n} \alpha}$: 这是译文 $\lwe_{k+1, m, q,\sqrt{5n} \alpha}$, 这是译文 $4 m \eps + \delta$. 这是译文 $\lwe_{k, m, q,\alpha}$ 这是译文 $\lwe_{k+1, m, q, \sqrt{5n} \alpha}$ (这是译文), 这是译文. 
\end{proof}

\begin{figure}[ht]%
\centering
\begin{tikzpicture}[node distance=1cm, auto]  
\tikzset{
    mynode/.style={rectangle,draw=black, top color=white, bottom color=yellow!5,thick, inner sep=0.5em, minimum height=2em, text centered},
    myarrow/.style={->, >=latex', shorten >=1pt, thick},
    mylabel/.style={text width=7em, text centered} 
}  
%\node[mynode] (final) {$\zolwe_{n, m',q, \sqrt{10n} \alpha}$};  
%\node[mynode, above=0.5cm of final] (atmost) {$\zolwe_{n, m,q, \le \sqrt{10n} \alpha}$}; 
\node[mynode] (atmost) {$\zolwe_{n, m,q, \le \sqrt{10n} \alpha}$}; 
\node[mynode, above right=1.2cm of atmost] (second) {$\lwe_{k+1, m, q, \sqrt{5n} \alpha}$};
\node[mynode, above left=1.2cm of atmost] (first) {$\extlwe^m_{k+1, n, q,\sqrt{5} \alpha,\{0,1\}^n}$};  
%\node[mynode, right=of second] (third) {discrete $\lwe^n_{k, m, q,\beta}$};
%\node[mynode, above=0.5cm of third] (thirdabove) {discrete $\lwe_{k, m, q,\beta}$};
%\node[mynode, above=0.5cm of thirdabove] (thirdaboveabove) {$\lwe_{k, m, q,\beta}$};
\node[mynode, above=0.5cm of first] (firstabove) {$\extlwe_{k+1, n, q, \sqrt{5} \alpha,\{0,1\}^n}$};
\node[mynode, above=0.5cm of firstabove] (firstaboveabove) {1$^{\mbox{\scriptsize{st}}}$ errorless $\lwe_{k+1, n,q,\alpha}$};
\node[mynode, above=5cm of atmost] (firstaboveaboveabove) {$\lwe_{k, m, q,\alpha}$};

%\draw[myarrow] (atmost)  --  (final);	
\draw[myarrow] (first)  --  (atmost);	
\draw[myarrow] (second)  --  (atmost);	
%\draw[myarrow] (third)  --  (atmost);	
%\draw[myarrow] (thirdabove)  --  (third);	
%\draw[myarrow] (thirdaboveabove)  --  (thirdabove);	
\draw[myarrow] (firstabove)  --  (first);	
\draw[myarrow] (firstaboveabove)  --  (firstabove);	
\draw[myarrow] (firstaboveaboveabove)  --  (firstaboveabove);	
\draw[myarrow] (firstaboveaboveabove)  --  (second);	
%\draw[myarrow] (firstaboveaboveabove)  --  (thirdaboveabove);	

%\node[mylabel, above left=0cm of final] (label1) {Lemma~\ref{lem:lelwe}};  
%\node[mylabel, below left=0cm of first] (label2) {Lemma~\ref{lem:zolwe}};  
\node[mylabel, above left=0cm of first] (label3) {Lemma~\ref{lem:multiinstanceextlwe}};  
\node[mylabel, above left=0cm of firstabove] (label4) {Lemma~\ref{lem:redtoextlwe}};  
\node[mylabel, below left=0cm of firstaboveaboveabove.west] (label5) {Lemma~\ref{lem:firsterrorlesshard}};  
%\node[mylabel, above right=0cm of third] (label6) {Lemma~\ref{lem:multiinstancelwe}};  
%\node[mylabel, above right=0cm of thirdabove] (label7) {Lemma~\ref{lem:conttodiscrete}};  
\node[rectangle,draw=black, thick, inner sep=0.5em, minimum width=5cm, fill=blue!10!white, fill opacity=0.8, text opacity=1, dashed, above=0.3cm of atmost] (emphasize) {Lemma~\ref{lem:zolwe}};
\end{tikzpicture} 
\caption{这是译文~\ref{thm:zolwe}}%
\label{fig:summary}%
\end{figure}




\subsection{这是译文}
\label{sec:firstiserrorless}

这是译文~\ref{lem:firsterrorlesshard} 这是译文. 

\begin{definition}\label{def:firsterrorless} 这是译文 $n,q \ge 1$ 这是译文 $\phi$ 这是译文 $\R$, 这是译文 $\lwe$ 这是译文 $\T_q^n \times \T_q^{}$ 这是译文 $\T_q^n \times \T$. 这是译文 $\vecs \in \{0,\ldots,q-1\}^n$, 这是译文 $A_{q,\vec{s},\{0\}}$ (这是译文 $\{0\}$ 这是译文) 这是译文 $A_{q,\vec{s},\phi}$.
\end{definition}


\begin{lemma}\label{lem:firsterrorlesshard} 这是译文 $n \ge 2$, $m, q \ge 1$, 这是译文 $\phi$, 这是译文 (这是译文) 这是译文 $\lwe_{n-1,m,q,\phi}$ 这是译文 $\lwe_{n,m,q,\phi}$ 这是译文 $\sum_p p^{-n}$, 这是译文 $q$.
\end{lemma}
这是译文 $q$ 这是译文 $q^{-n}$. 这是译文 $q$ 这是译文 \[
\sum_{k \ge 2} k^{-n} \le 2^{-n} + \int_{2}^\infty t^{-n} \mathrm{d} t \le 2^{-n+2},
\] 这是译文 $n$ 这是译文.
\begin{proof}
这是译文 $\veca'$ 这是译文  $\{0,\ldots,q-1\}^n$. 这是译文 $r$ 这是译文 $\veca'$.  这是译文 $q$, 这是译文  \[
 \sum_{p~\text{prime}, \ p|q} p^{-n}.
\] 这是译文 $\matU \in \Z^{n \times n}$ 这是译文 $q$ 这是译文 $\veca'$. 这是译文 $n\times n$ 这是译文 $\matR$ 这是译文 $\matR \veca' = (r,0,\ldots,0)^T$. 这是译文 $\matR^{-1} \cdot \mathrm{diag}(r,1,\ldots,1)$ 这是译文 $s_0 \in \{0,\ldots,q-1\}$. 这是译文 $(\veca'/q, s_0/q)$. 这是译文 $(\veca,b)$ 这是译文 $d \in \T_q$, 这是译文 $(\matU(d|\veca),b+(s_0 \cdot d))$ 这是译文 (这是译文 $b$ 这是译文 $\T_q$), 这是译文 $2^{-n+1}$; 这是译文 $A_{q,\vec{s},\phi}$, 这是译文 $A_{q,\vecs',\{0\}}$ 这是译文 $A_{q,\vecs',\phi}$, 这是译文 $2^{-n+1}$, 这是译文 $\vecs' = (\matU^{-1})^T (s_0|\vecs) \bmod q$. 这是译文 $\matU$, 这是译文 $q$, 这是译文 $\Z_q^n$, 这是译文 $\vecs'$ 这是译文 $\{0,\ldots,q-1\}^n$.
\end{proof}

\subsection{这是译文}
\label{sec:extlwe}

这是译文 $\extlwe$. (这是译文~\cite{DBLP:conf/pkc/Alperin-SheriffP12}, 这是译文.)

\begin{definition}\label{def:extlwe} 这是译文 $n,m,q,t \ge 1$, $\calZ \subseteq \Z^m$, 这是译文 $\chi$ 这是译文 $\frac{1}{q}\Z^m$, 这是译文 $\extlwe^t_{n,m,q,\chi,\calZ}$ 这是译文 $\vecz \in \calZ$ 这是译文 \[
(\matA,(\vecb_i)_{i \in [t]},(\inner{\vece_i,\vecz})_{i \in [t]}) \in \T_q^{n \times m} \times (\T_q^m)^t \times ({\textstyle \frac 1 q} \Z)^t.
\] 这是译文, $\matA \in \T_q^{n \times m}$ 这是译文, $\vece_i \in {\textstyle \frac 1 q}\Z^m$ 这是译文 $\chi$, 这是译文 $\vecb_i = \matA^T \vecs_i + \vece_i \bmod 1$ 这是译文 $\vecs_i \in \{0,\ldots,q-1\}^n$ 这是译文 $\vecb_i$ 这是译文 $\T_q^m$ 这是译文. 
\end{definition}


这是译文 $t=1$, 这是译文 $t$. 这是译文 $\chi$ 这是译文 $D_{q^{-1}\Z^{m},\alpha}$ 这是译文 $\alpha>0$, 这是译文 $\chi$ 这是译文 $\alpha$.  这是译文 $\lwe$ 这是译文 $\extlwe$ 这是译文 $\cZ = \{ 0^m \}$.  这是译文 $\calZ$.



\begin{definition}\label{def:qualityofhintset} 这是译文 $\xi>0$ 这是译文 $\calZ \subseteq \Z^m$ 这是译文 $\calZ$ 这是译文 \emph{这是译文} $\xi$ 这是译文 $\vecz \in \calZ$, 这是译文 $\matU \in \Z^{m \times m}$ 这是译文 $\matU'\in \Z^{m \times (m-1)}$ 这是译文 $\matU$ 这是译文 $\matU'$ 这是译文 $\vecz$ 这是译文~$\xi$.
\end{definition}

这是译文 $\matU'$  这是译文 $\vecz$, 这是译文 $\{\vecb \in \Z^m:
  \inner{\vecb,\vecz} = 0\}$. 这是译文.

\begin{claim}\label{clm:qualityofzeroone} 这是译文 $\calZ=\{0,1\}^m$ 这是译文 $2$.
\end{claim}
\begin{proof}
这是译文 $\vecz \in \calZ$ 这是译文 $k \ge 1$ 这是译文 $1$ 这是译文 $m-k$ 这是译文 $0$. 这是译文 $\matU$ 这是译文 $1$这是译文 $(-1,\ldots,-1,0,\ldots,0)$ 这是译文 $-1$ 这是译文 $k-1$ 这是译文 $\vecz$. 这是译文 $\matU$ 这是译文 $1$ 这是译文 $1$.
\end{proof}

\onote{todo: add more quality bounds such as vectors of bounded norm}

\begin{lemma}\label{lem:redtoextlwe} 这是译文 $\calZ \subseteq \Z^m$ 这是译文 $\xi>0$. 这是译文 $n,q \ge 1$, $\eps \in (0,1/2)$, 这是译文 $\alpha,r  \ge (\ln(2m(1+1/\eps))/\pi)^{1/2}/q$, 这是译文 (这是译文) 这是译文 $\lwe_{n,m,q,\alpha}$ 这是译文 $\extlwe_{n,m,q,(\alpha^2 \xi^2 + r^2)^{1/2},\calZ}$ 这是译文 $33 \eps/2$. 
\end{lemma}
\begin{proof}
这是译文 $\vecz \in \calZ$. 这是译文 $\matU \in \Z^{m \times m}$ 这是译文 $\vecz$ 这是译文~\ref{def:qualityofhintset}, 这是译文 $\matU' \in \Z^{m \times (m-1)}$ 这是译文 $\matU$. 这是译文 $m$ 这是译文 $(\matA, \vecb) \in \T_q^{n \times m} \times (\T_q^{} \times \T^{m-1})$. 这是译文 $\vecf$ 这是译文 $m$-这是译文 $D_{\alpha(\xi^2 \matI - \matU' \matU'^T)^{1/2}}$,  这是译文 $\xi^2 \matI - \matU' \matU'^T$ 这是译文 $\matU$. 这是译文 \begin{equation}\label{eq:outputofextlwereduction}
(\matA'= \matA \matU^T,
  \vecb' + \vecc,
	\inner{\vecz , \vecf + \vecc})
	\in \T_q^{n \times m} \times \T_q^m \times {\textstyle{\frac{1}{q}}}\Z,
\end{equation} 这是译文 $\vecb'=\matU \vecb + \vecf$, 这是译文 $\vecc$ 这是译文 $D_{q^{-1}\Z^m - \vecb',r}$ (这是译文~\ref{lem:gpv}). 

这是译文., $\matA$ 这是译文~$\T_q^{n \times m}$ 这是译文 $\vecb = \matA^T \vecs + \vece \in \T^m$ 这是译文 $\vecs \in \{0,\ldots,q-1\}^n$ 这是译文 $\vece \in \R^m$ 这是译文~$0$,  这是译文 $m-1$ 这是译文~$D_\alpha$. 这是译文 $\matU$ 这是译文, $\matA' = \matA \matU^T$ 这是译文 $\T_q^{n \times m}$ 这是译文 $\matA'$ 这是译文~\eqref{eq:outputofextlwereduction}. 这是译文, \[
\vecb' =
\matU \vecb + \vecf =
\matA'^T \vecs + \matU \vece + \vecf.
\] 这是译文 $\matU \vece$ 这是译文 $D_{\alpha \matU'}$, 这是译文~$\matU \vece + \vecf$ 这是译文 \emph{这是译文} 这是译文~$D_{\alpha \xi}$. 这是译文 $\matA'^T \vecs \in \T_q^m$, 这是译文 $q^{-1}\Z^m - \vecb'$ 这是译文 $q^{-1}\Z^m - (\matU \vece + \vecf)$, 这是译文 $\vecc$ 这是译文 $D_{q^{-1}\Z^m - (\matU \vece + \vecf),r}$. 这是译文~\ref{lem:peikertdiscretegauss} 这是译文 $r \ge \eta_\eps(q^{-1}\Z^m)$  这是译文~\ref{lem:boundonsmoothing}, 这是译文 $\matU \vece + \vecf + \vecc$ 这是译文 $8 \eps$  这是译文 $D_{q^{-1}\Z^m,(\alpha^2 \xi^2 + r^2)^{1/2}}$. 这是译文~\eqref{eq:outputofextlwereduction} 这是译文 $\matU$ 这是译文 $\vece$ 这是译文, \[
\inner{\vecz, \vecf + \vecc} = \inner{\vecz, \matU\vece + \vecf + \vecc},
\] 这是译文 $\vecz$, 这是译文.

这是译文 $\matA$ 这是译文~$\T_q^{n\times m}$ 这是译文 $\vecb$ 这是译文 $\T_q \times \T^{m-1}$. 这是译文~\ref{lem:smoothingcontinuous}, 这是译文 $\alpha \ge \eta_{\eps/m}(q^{-1}\Z)$ (这是译文~\ref{lem:boundonsmoothing}), 这是译文~$(\matA,\vecb)$ 这是译文 $\eps/2$ 这是译文 $(\matA, \vece' + \vece)$ 这是译文 $\vece'$ 这是译文~$\T_q^m$, 这是译文 $\vece$ 这是译文 $m-1$ 这是译文~$D_\alpha$. 这是译文 $(\matA,\vece'+\vece)$. 这是译文~\eqref{eq:outputofextlwereduction} 这是译文 $\T_q^{n \times m}$ 这是译文 $\vecb' = \matU \vece' + \matU \vece + \vecf$ 这是译文 $\matU \vece' \in \T_q^m$, 这是译文 $q^{-1}\Z^m - \vecb'$ 这是译文 $q^{-1}\Z^m - (\matU \vece + \vecf)$, 这是译文 $\vecc$ 这是译文 $\vece'$. 这是译文~\eqref{eq:outputofextlwereduction} 这是译文 $\vece'$ 这是译文 $\matU \vece'$ 这是译文 $\T_q^m$ (这是译文 $\matU$ 这是译文), 这是译文.
\end{proof}

这是译文 ($t \ge 1$) 这是译文. 

%The proof is by a straightforward hybrid argument and is omitted.

\onote{removed $\calD$ from the following lemma since we never  use any distribution on secrets but the uniform one}

\begin{lemma}\label{lem:multiinstanceextlwe} 这是译文 $n, m, q, \chi, \cZ$ 这是译文~\ref{def:extlwe} 这是译文 $\chi$ 这是译文 $t \ge 1$ 这是译文 (这是译文) 这是译文 $\extlwe_{n,m,q,\chi,\cZ}$ 这是译文 $\extlwe^t_{n,m,q,\chi,\cZ}$ 这是译文~$t$.
%\dnote{I changed~$\phi$ to $\chi$, for consistency with 
%definition~\ref{def:extlwe}.} 
\end{lemma}

这是译文 $\vc{s}$ 这是译文 (这是译文) 这是译文 $\vc{s}$.

\begin{proof}
这是译文 $\cA$ 这是译文 $\extlwe^t_{n,m,q,\chi,\cZ}$, 这是译文 $\vc{z}$ 这是译文 $\cA$ 这是译文 (这是译文) 这是译文 $H_i$ 这是译文 \[
\left(\mx{A}, \{\vc{b}_1, \ldots, \vc{b}_i, \vc{u}_{i+1}, \dots, \vc{u}_t\}, \vc{z}, \{\langle\vc{z}, \vc{e}_i\rangle\}_{i\in[t]} \right)~,\]  这是译文~$\vc{u}_{i+1}, \dots, \vc{u}_t$ 这是译文~$\T_q^m$. 这是译文 $\adv[\cA] = \abs{\Pr[\cA(H_0)]- \Pr[\cA(H_t)]}$.

这是译文 $\cB$ 这是译文 $\extlwe_{n,m,q,\chi,\cZ}$: 这是译文, $\cB$ 这是译文 $\cA$ 这是译文 $\vc{z}$ 这是译文 $\vc{z}$. 这是译文 $(\mx{A}, \vc{d}, \vc{z}, y)$ 这是译文 $\extlwe_{n,m,q,\chi,\cZ}$, 这是译文 $\cB$ 这是译文 $i^* \getsr [t]$, 这是译文 $\vc{s}_1, \ldots, \vc{s}_{i^*-1}\getsr \bbZ^n_q$, $\vc{e}_1, \ldots, \vc{e}_{i^*-1}, \vc{e}_{i^*+1}, \ldots, \vc{e}_t \getsr \chi^m$\onote{replaced $\phi$ with $\chi$}, $\vc{u}_{i^*+1}, \ldots, \vc{u}_t \getsr \bbT_q^m$. 这是译文 $\vc{b}_i = \mx{A}^T\cdot\vc{s}_i+\vc{e}_i\pmod{1}$, 这是译文 $\cA$: \[
\left(\mx{A}, \{ \vc{b}_1, \ldots, \vc{b}_{i^*-1}, \vc{d}, \vc{u}_{i^*+1}, \ldots, \vc{u}_t\}, \vc{z}, \{\langle\vc{z},\vc{e}_1\rangle, \ldots,\langle\vc{z}, \vc{e}_{i^*-1}\rangle, y, \langle\vc{z}, \vc{e}_{i^*+1}\rangle, \ldots, \langle\vc{z}, \vc{e}_t\rangle  \} \right)~.
\] 这是译文, $\cB$ 这是译文 $\cA$ 这是译文.

这是译文 $\cB$ 这是译文 $P_0=(\mx{A}, \vc{b}, \vc{z}, \vc{z}^T \cdot \vc{e})$ 这是译文~$\vc{b} = \mx{A}^T\cdot\vc{s}+\vc{e}\pmod{1}$, 这是译文 $\cB$ 这是译文 $\cA$ 这是译文 $H_{i^*}$. 这是译文 $\cB$ 这是译文 $P_1 = (\mx{A}, \vc{u}, \vc{z}, \vc{z}^T \cdot \vc{e})$ 这是译文~$\vc{u} \getsr \bbT_q^m$, 这是译文 $\cB$ 这是译文 $\cA$ 这是译文~$H_{i^*-1}$.

这是译文 $i^*$ 这是译文 $[t]$, 这是译文 \begin{eqnarray*}
t \, \adv[\cB] & = & t \, \abs{\Pr[\cB(P_0)] - \Pr[\cB(P_1)]}\\
  & = &
	\abs[\bigg]{{\sum_{i^*\in[t]} \Pr[\cA(H_{i^*})] - \sum_{i^*\in[t]} \Pr[\cA(H_{i^*-1})]}} \\
	& = &
	\abs{\Pr[\cA(H_{t})]- \Pr[\cA(H_{0})] } \\
	& = & \adv[\cA]~,
\end{eqnarray*} 这是译文.
\end{proof}




%%% Local Variables:
%%% mode: latex
%%% TeX-master: "lwehard"
%%% End:


\subsection{这是译文}\label{sec:theshortsecretreduction}

\begin{lemma}\label{lem:zolwe} 这是译文 $k,n,m,q \in \bbN$, $\epsilon \in (0,1/2)$, 这是译文 $\delta,\alpha,\beta,\gamma>0$ 这是译文 $n \ge k \log_2 q + 2 \log_2 (1/\delta)$, $\beta \ge \sqrt{2 \ln(2n(1+1/\eps))/\pi}/q$, $\alpha = \sqrt{2n} \beta$, $\gamma = \sqrt{n}\beta$.  这是译文 (这是译文) 这是译文 $\zolwe_{n,m,q,\le\alpha}$ 这是译文  $\extlwe^m_{k,n,q,\beta,\{0,1\}^n}$, $\lwe_{k,m,q,\gamma}$, 这是译文  $\extlwe^m_{k,n,q,\beta,\{0^n\}}$, 这是译文 $\cB_1$, $\cB_2$, 这是译文 $\cB_3$ 这是译文 (这是译文)  这是译文 $\cA$, 这是译文 \[
\adv[\cA] \le \adv[\cB_1]+\adv[\cB_2]+\adv[\cB_3]+4m\epsilon + \delta \,.
\]
\end{lemma}

这是译文 (这是译文) 这是译文 $\extlwe^m_{k,n,q,\beta,\{0,1\}^n}$ 这是译文 $\extlwe^m_{k,n,q,\beta,\{0^n\}}$, 这是译文 $\zolwe_{n,m,q,\le\alpha}$ 这是译文 $\extlwe^m_{k,n,q,\beta,\{0,1\}^n}$ 这是译文 $\lwe_{k,m,q,\gamma}$.

这是译文 $k \log_2 q + 2 \log_2 (1/\delta)$, 这是译文 $\zolwe$.


\def\veh{\hat{\vc{e}}}


\begin{proof}
这是译文 $k,n,m,q,\eps, \alpha, \beta, \gamma$ 这是译文 $\vc{z} \getsr \binset^n$ 这是译文 $\vc{e} \getsr D_{\alpha'}^m$ 这是译文 $\alpha' = \sqrt{\beta^2\|\vc{z}\|^2+ \gamma^2} \le \sqrt{2n}\beta=\alpha$. 这是译文 $\mx{A}\getsr \bbT_q^{n \times m}$, $\vc{u} \getsr \bbT^m$, 这是译文 $\vc{b} \getsd \mx{A}^T\cdot\vc{z}+\vc{e}\pmod{1}$. 这是译文~$\cA$ 这是译文 $(\mx{A}, \vc{b})$ 这是译文 $(\mx{A}, \vc{u})$.
%\dnote{I modified the last sentence a bit}

这是译文 $H_0$ 这是译文 $(\mx{A}, \vc{b})$ 这是译文 $H_1$ 这是译文  \[ H_1 = (\mx{A}, \mx{A}^T\vc{z}-\mx{N}^T\vc{z}+\veh \bmod 1 ),\] 这是译文 $\mx{N}\getsr D^{n \times m}_{q^{-1}\bbZ, \beta}$ 这是译文 $\veh \getsr D_{\gamma}^m$. 这是译文 $\| \vc{z} \| \le \sqrt{n}$ 这是译文 $\beta \ge \sqrt{2} \eta_\eps(\Z^n)/q$ (这是译文~\ref{lem:boundonsmoothing}), 这是译文~\ref{lem:smoothip} 这是译文 $-\mx{N}^T\vc{z}+\veh$ 这是译文 $D^m_{\alpha'}$ 这是译文 $4 m \epsilon$. 这是译文 \begin{equation}\label{eq:h0}
\abs{\Pr[\cA(H_0)] - \Pr[\cA(H_1)]} \le 4m\epsilon~.
\end{equation}


这是译文 $H_2$ 这是译文 $\mx{B} \getsr \bbT_q^{k \times m}$ 这是译文 $\mx{C} \getsr \bbT_q^{k \times n}$. 这是译文 $\mAh \getsd q\mx{C}^T\cdot \mx{B} + \mx{N}\pmod{1}$. 这是译文, \[
H_2 = (\mAh, \mAh^T\cdot\vc{z} - \mx{N}^T\vc{z}+\veh) = (\mAh, q\mx{B}^T\cdot \mx{C}\cdot\vc{z}+\veh)~.
\]

这是译文 $\cB_1$ 这是译文 $\extlwe^m_{k,n,q,\beta,\{0,1\}^n}$, 这是译文 \begin{equation}\label{eq:h1}
\adv[\cB_1] = \abs{\Pr[\cA(H_1)]-\Pr[\cA(H_2)]}~.
\end{equation} 这是译文 $H_1, H_2$ 这是译文 $(\mx{C}, \mx{A}, \mx{N}^T\vc{z})$ 这是译文 $(\mx{C}, \mAh, \mx{N}^T\vc{z})$ 这是译文 $\extlwe^m_{k,n,q,\beta,\{0,1\}^n}$ 这是译文 (这是译文 $q\cdot\mx{B}$ 这是译文 $m$ 这是译文), 这是译文 $\cB_1$ 这是译文~$\cA$.

这是译文 $H_3 = (\mAh, \mx{B}^T\cdot \vc{s}+\veh)$, 这是译文 $\vc{s} \getsr \bbZ_q^k$. 这是译文 \begin{equation}\label{eq:h2}
\abs{\Pr[\cA(H_2)]- \Pr[\cA(H_3)]} \le \delta
\end{equation} 这是译文 (这是译文~\ref{lem:lhl}), 这是译文 $H_2, H_3$ 这是译文 $(\mx{C}, q\mx{C}\cdot\vc{z})$ 这是译文 $(\mx{C}, \vc{s})$ 这是译文 $\delta$.

这是译文: $H_4=(\mAh, \vc{u})$. 这是译文 $\cB_2$ 这是译文 $\lwe_{k,m,q,\gamma}$ 这是译文 \begin{equation}\label{eq:h3}
\adv[\cB_2] = \abs{\Pr[\cA(H_3)]- \Pr[\cA(H_4)]}~,
\end{equation} 这是译文 $H_3, H_4$ 这是译文 $(\mx{B}, \mx{B}^T\vc{s}+\veh), (\mx{B}, \vc{u})$.

这是译文 $\mAh$ 这是译文: $H_5 = (\mx{A}, \vc{u})$. 这是译文 $\cB_3$ 这是译文 $\extlwe^m_{k,n,q,\beta,\{0^n\}}$ 这是译文 \begin{equation}\label{eq:h4}
\adv[\cB_3] = \abs{\Pr[\cA(H_4)]- \Pr[\cA(H_5)]}~.
\end{equation} 这是译文.~\eqref{eq:h4} 这是译文.~\eqref{eq:h1}: 这是译文 $H_4$, $H_5$ 这是译文 $(\mx{C}, \mAh)$ 这是译文 $(\mx{C}, \mx{A})$ 这是译文 $\extlwe^m_{k,n,q,\beta,\{0^n\}}$ 这是译文 (这是译文 $q\cdot\mx{B}$ 这是译文 $m$ 这是译文), 这是译文 $\cB_3$ 这是译文~$\cA$.

这是译文.~\eqref{eq:h0},~\eqref{eq:h1},~\eqref{eq:h2},~\eqref{eq:h3},~\eqref{eq:h4}, 这是译文.
\end{proof}












%%% Local Variables:
%%% mode: latex
%%% TeX-master: "lwehard"
%%% End:



\section{这是译文}
\label{sec:modexpansion}



这是译文~\ref{cor:mod-reduction} 这是译文~\ref{cor:mod-dim-tradeoff} 这是译文~$\calD$ 这是译文~$\Z^{n}$ 这是译文 $(B,\delta)$-这是译文~$B,\delta\geq 0$ 这是译文~$\vec{x} \gets
\calD$ 这是译文 $B$ 这是译文 $\delta$. 

\begin{theorem}
  \label{thm:mod-dim-switch}  这是译文 $m,n, n',q,q' \geq 1$ 这是译文  $\matG \in \Z^{n' \times n}$ 这是译文 $\Lambda =
  \frac{1}{q'} \matG^{T} \Z^{n'} + \Z^{n}$ 这是译文 $\matB$,  这是译文 $\calD$ 这是译文 $(B,\delta)$-这是译文  $\Z^{n}$.  这是译文 $\alpha, \beta > 0$ 这是译文~$\eps \in (0,1/2)$ 这是译文  \[\beta^{2} \geq
  \alpha^{2} + (4/\pi) \ln(2n(1+1/\eps)) \cdot (\max\set{q^{-1}, \length{\gs{\matB}}} \cdot B)^{2}.\]  这是译文 (这是译文) 这是译文~$\lwe_{n,m,q,\le \alpha}(\calD)$ 这是译文 $\lwe_{n',m,q',\leq \beta}(\matG
  \cdot \calD)$ 这是译文 $\delta + 14 \eps m$.
\end{theorem}

这是译文 $\|\widetilde \matB\|$ 这是译文~\ref{lem:gpv}. 这是译文 $\Z_{q'}^{n'}$, 这是译文~\ref{def:lwe}. 这是译文 $\lwe_{n',q',\leq \beta}$  这是译文~$\lwe_{n',q',\beta}$ (这是译文~\ref{lem:lelwe}).

这是译文 (这是译文 \lwe 这是译文), 这是译文 $n'=n$, $\matG=\matI$ 这是译文 $n$-这是译文 $\matB = \matI/q'$.  这是译文 $q \ge q' \ge \sqrt{2\ln(2n(1+1/\eps))} \cdot
(B/\alpha)$ 这是译文 $\beta=\sqrt{2} \alpha$, 这是译文~$B/\alpha$, 这是译文~$\alpha$ 这是译文.

\begin{corollary}
  \label{cor:mod-reduction}  这是译文 $m,n \geq 1$, $q \geq q' \geq 1$, $(B,\delta)$-这是译文~$\calD$ 这是译文~$\Z^{n}$, $\alpha,\beta>0$ 这是译文~$\eps \in (0,1/2)$ 这是译文  \[
	\beta^{2} \geq \alpha^{2} + (4/\pi) \ln(2n(1+1/\eps))
  \cdot (B/q')^2,
	\]  这是译文  $\lwe_{n,m,q,\le \alpha}(\calD)$ 这是译文  $\lwe_{n,m,q',\le \beta}(\calD)$ 这是译文 $\delta + 14 \eps m$.
\end{corollary}

这是译文 \lwe (这是译文~\ref{lem:normalform}), 这是译文 $\calD = D_{\Z^{n}, \sqrt{2} \alpha q}$, 这是译文 (这是译文~\ref{thm:searchtodecisionmicpei}), 这是译文 (这是译文) $\lwe$ 这是译文 \emph{这是译文} 这是译文 $q$. 这是译文 $\calD = D_{\Z^{n}, r}$ 这是译文 $(C r \sqrt{n \log(n/\delta)}, \delta)$-这是译文 $C>0$, 这是译文 $n$ 这是译文. (这是译文~$(r\sqrt{n},2^{-n})$-这是译文~\cite[Lemma~1.5]{banaszczyk93:_new}, 这是译文 $n$.)

\begin{corollary}
  \label{cor:mod-reduction-2}  这是译文 $\delta \in (0,1/2)$, $m \ge n \geq 1$, $q' \ge 25$. 这是译文 $q \in [q',2q')$ 这是译文 $2$ 这是译文 $q'$  这是译文 $\alpha \ge \sqrt{\ln(2n(1+16/\delta)/\pi)}/q$.  这是译文 (这是译文) 这是译文 $\lwe_{n,m,q, \alpha}$ 这是译文 $\lwe_{n,m',q', \le \beta}$  这是译文 $m'=m-(16 n + 4 \ln\ln q)$ 这是译文  \[
	\beta = C \alpha \sqrt{n} \sqrt{\log(n/\delta)\log(m/\delta)}
	\]  这是译文 $C>0$, 这是译文 $\zeta$ 这是译文 $(\zeta-\delta)/4$.
\end{corollary}


  这是译文  $n=kn'$ 这是译文 $k \ge 1$, 这是译文~$q' = q^{k}$.  这是译文~$\vec{G} =
  \vec{I}_{n'} \otimes \vec{g}$, 这是译文~$\vec{g} = (1,q,q^2, \ldots,
  q^{k-1})^T \in \Z^k$.  这是译文 $\Lambda = q^{-k} \vec{G}^T
  \Z^{n'} + \Z^n$.  这是译文~$\Lambda$ 这是译文 \[
\vec{B}= \vec{I}_{n'} \otimes 
\begin{bmatrix}
q^{-1} & q^{-2} & \cdots & q^{-k} \\
       & q^{-1} & \cdots & q^{1-k} \\
       &        & \ddots & \vdots \\
       &		&		 & q^{-1}
     \end{bmatrix} \in \R^{n \times n};
\] 这是译文~$\vec{B}$ 这是译文~$\Lambda$ 这是译文 $\gs{\matB} = q^{-1} \matI$ 这是译文 $\length{\gs{\matB}}=q^{-1}$.  这是译文~$n \log q = n' \log q'$ 这是译文~$\calD = D_{\Z^{n}, \alpha q}$ (这是译文~\ref{lem:normalform}), 这是译文 $(\alpha q
\sqrt{n}, 2^{-n})$-这是译文~$\sqrt{n}$ 这是译文. 

\begin{corollary}
  \label{cor:mod-dim-tradeoff}  这是译文 $n, m,q \ge 1$, $k \ge 1$ 这是译文~$n$,  $(B,\delta)$-这是译文~$\calD$ 这是译文~$\Z^{n}$,  $\alpha,\beta>0$, 这是译文~$\eps \in (0,1/2)$ 这是译文 \[
\beta^{2} \geq \alpha^{2} + (4/\pi)
  \ln(2n(1+1/\eps)) \cdot (B/q)^{2},
\]  这是译文~$\lwe_{n,m,q,\le \alpha}(\calD)$  这是译文~$\lwe_{n/k,m,q^k,\le \beta}(\vec{G} \cdot \calD)$ 这是译文 $\delta+14 \eps m$,  这是译文~$\vec{G} = \vec{I}_{n/k} \otimes  (1,q,q^2, \ldots, q^{k-1})^T$.
\end{corollary}

\noindent 这是译文~\ref{thm:mod-dim-switch} 这是译文.


\begin{lemma}
  \label{lem:reverse-main}  这是译文~\ref{thm:mod-dim-switch}, 这是译文  \[
	r
  \geq \max\set{q^{-1}, \length{\gs{\matB}}} \cdot
  \sqrt{2\ln(2n(1+1/\epsilon))/\pi}. 
	\] 这是译文 $\T_{q}^{n} \times \T$ 这是译文 $\T_{q'}^{n'} \times \T$, 这是译文:
  \begin{itemize}[itemsep=0pt]
  \item 这是译文 $4 \epsilon$ 这是译文.
  \item 这是译文  $A_{q,\vecs,D_{\alpha}}$ 这是译文 $\vecs \in \Z^{n}$ 这是译文  $\length{\vecs} \leq B$, 这是译文 $10 \epsilon$ 这是译文 $A_{q',\matG\vecs,D_{\alpha'}}$, 这是译文  $(\alpha')^{2} = \alpha^{2} + r^{2}(\length{\vecs}^{2}+B^{2}) \leq
    \alpha^{2} + 2(rB)^{2}$. 
%\dnote{I think we use the condition $\length{\vecs} \leq B$ only for that last inequality. 
%If this is indeed the case, I would be in favour of removing the condition, and commenting at the end of this item that 
%if~$\length{\vecs} \leq B$, then $\alpha^{2} + r^{2}(\length{\vecs}^{2}+B^{2}) \leq
%    \alpha^{2} + 2(rB)^{2}$}\onote{don't we also need it when we apply Lemma~\ref{lem:smoothip} towards the end of the proof? (I didn't check too carefully)}
%\dnote{Why? we don't need $\vec{s}$ small too apply it, do we?}\onote{yes, I think we do; look carefully at the lemma...}
  \end{itemize}
\end{lemma}

\begin{proof}
  这是译文 $\T_{q}^{n}$ 这是译文  $\T_{q'}^{n'}$, 这是译文-$1$ 这是译文 $\T_{q}^{n}$ 这是译文 $\vecs \in \Z^{n}$, 这是译文 $\T_{q'}^{n'}$ 这是译文 $\matG \vecs \in \Z^{n'}$, 这是译文 $\veca \in \T_{q}^{n}$ (这是译文) 这是译文 $\veca' \in
  \T_{q'}^{n'}$, 这是译文  \[ \inner{\veca', \matG \vecs} = \inner{\matG^{T} \veca', \vecs}
  \approx \inner{\veca,\vecs} \bmod 1 \] 这是译文 (这是译文) $\vecs \in
  \Z^{n}$.  这是译文 $\veca'$ 这是译文  $\matG^{T} \veca' \approx \veca \bmod \Z^{n}$, 这是译文 $r$.
  

  % In order to formally describe how the reduction samples $\veca'$, we
  % first define an $n$-dimensional lattice \[ \Lambda = q'^{-1} \cdot
  % \lat(\matG^{T}) + \Z^{n} = q'^{-1} \cdot \lat(\matI_{n'} \otimes
  % \vecg) + \Z^{n}, \] and claim that it has as a basis the
  % upper-triangular matrix \[ \matB =
  % \matI_{n'} \otimes
  % \begin{bmatrix}
  %   (q'_{1})^{-1} & \star & \cdots & \star \\
  %   & (q'_{2})^{-1} & \cdots & \star \\
  %   & & \ddots & \vdots \\
  %   & & & (q'_{k})^{-1}
  % \end{bmatrix} \in \R^{n \times n}, \] where the rightmost column of
  % the above displayed $k$-by-$k$ matrix is $q'^{-1} \vecg$, and the
  % $j$th column is obtained by multiplying the $(j+1)$st column by
  % $q'_{j+1}$ and reducing modulo $1$, for $j=k-1,\ldots, 1$.  Indeed,
  % $\matB$ is a basis of $\Lambda$ because all of its columns are in
  % $\Lambda$ by construction, and because
  % $\det(\matB)=q'^{-n'}=\abs{\Lambda/\Z^{n}}^{-n'}=\det(\Lambda)$.  Also
  % notice that because $\matB$ is upper triangular, its Gram-Schmidt
  % orthogonalization (from left to right) is just the diagonal matrix
  % with $(q'_{i})^{-1}$ in the $i$th diagonal, so $\max_{i}
  % \length{\gs{\vecb_{i}}} \leq q_{\min}^{-1}$ and $r \geq
  % \sqrt{2}\smootheps(\Lambda)$ by~\cite[Theorem
  % 3.1]{DBLP:conf/stoc/GentryPV08}\cnote{Put this in prelims...}.
  
  这是译文 $(\veca, b) \in \T_{q}^{n} \times \T$, 这是译文:
  \begin{itemize}
  \item 这是译文 $\vecf \gets D_{\Lambda-\veca,r}$ 这是译文~\ref{lem:gpv} 这是译文  $\matB$, 这是译文 $\vecv = \veca + \vecf \in
    \Lambda/\Z^{n}$. (这是译文 $\Lambda-\veca$ 这是译文 $\veca =
      \bar{\veca}+\Z^{n}$ 这是译文 $\Z^{n} \subseteq \Lambda$.) 这是译文 $\veca' \in \T_{q'}^{n'}$ 这是译文  $\matG^{T} \veca' = \vecv \bmod \Z^{n}$. 这是译文 $\matG^{T}
      \veca' = \veczero \bmod \Z^{n}$, 这是译文  $\matG^{T} \veca' = \vecv \bmod \Z^{n}$.
  \item 这是译文 $e' \gets D_{rB}$ 这是译文 $b' = b + e' \in \T$.  
  \item 这是译文 $(\veca',b')$.
  \end{itemize}

  这是译文 $\veca'$ 这是译文  (这是译文) 这是译文 $b$ 这是译文~$e'$ 这是译文~$\veca'$, 这是译文 $b \in
  \T$ 这是译文 $\vecv \in \Lambda/\Z^{n}$ 这是译文 (这是译文) 这是译文 $\vecv$ 这是译文  $\veca'$ 这是译文 $\matG^{T} \veca' = \vecv \bmod \Z^{n}$.  这是译文 $\bar{\veca} \in \T_q^n$ 这是译文 $\bar{\vecf} \in \Lambda - \bar{\veca}$, 这是译文~\ref{lem:smoothing}  (这是译文 $r \ge \smootheps(\Lambda)$ 这是译文~\ref{lem:boundonsmoothing}) 这是译文  \begin{align}\label{eq:roundingonelatticeanother}
	  \Pr[\veca = \bar{\veca} \wedge \vecf = \bar{\vecf}] &= q^{-n} \cdot \rho_r(\bar{\vecf}) / \rho_r(\Lambda-\bar{\veca}) \nonumber \\
		& \in C \bracks[\big]{1,\tfrac{1+\epsilon}{1-\epsilon}}\cdot \rho_r(\bar \vecf).
	\end{align}  这是译文 $C=q^{-n}/\rho_{r}(\Lambda)$ 这是译文 $\bar{\veca}$ 这是译文 $\bar{\vecf}$.  这是译文 $\bar \veca, \bar \vecf$ 这是译文 $\bar \veca + \bar \vecf = \bar \vecv$, 这是译文 $\bar{\vecv} \in \Lambda/\Z^{n}$,  \begin{align*}
    \Pr[\vecv = \bar{\vecv}] 
    &\in C \bracks[\big]{1,\tfrac{1+\epsilon}{1-\epsilon}} \cdot
    \rho_{r}(q^{-1} \Z^{n} + \bar{\vecv}).
  \end{align*}  这是译文 $r
  \geq \smootheps(q^{-1} \Z^{n})$ (这是译文~\ref{lem:boundonsmoothing}), 这是译文~\ref{lem:smoothing}  这是译文~$\Pr[\vecv = \bar{\vecv}] \in \bracks[\big]{\tfrac{1-\epsilon}{1+\epsilon},\tfrac{1+\epsilon}{1-\epsilon}} C'$ 这是译文~$C'$ 这是译文~$\bar{\vec{v}}$. 这是译文~\ref{clm:separationdist}, 这是译文~$\veca'$ 这是译文~$1-((1-\eps)/(1+\eps))^2 \le 4\epsilon$ 这是译文.

  这是译文 $A_{q,\vecs,D_{\alpha}}$  这是译文 $A_{q', \matG \vecs, D_{\beta}}$.  这是译文 $(\veca,b=\inner{\veca,\vecs}+e)$, 这是译文 $e
  \gets D_{\alpha}$.  这是译文 $\veca'$ 这是译文 (这是译文)  这是译文~$\T_{q'}^{n'}$.  这是译文 $\overline{\veca'} \in \T_{q'}^{n'}$ 这是译文 $\veca'$,  这是译文 $\bar \vecv = \matG^{T} \overline{\veca'} \bmod \Z^{n}$.  这是译文  \[ b' = \inner{\veca,\vecs} + e + e' = \inner{\overline{\veca'}, \matG \vecs} +
  e + \inner{-\vecf, \vecs} + e' \bmod 1. \]  这是译文~\ref{clm:separationdist} 这是译文~\eqref{eq:roundingonelatticeanother} (这是译文 $\vecf=\bar \vecf$ 这是译文 $\veca = \bar \vecv - \bar \vecf \bmod \Z^n$),  这是译文 $-\vecf$ 这是译文~$1-(1-\eps)/(1+\eps) \le 2 \eps$ 这是译文  $D_{q^{-1}\Z^{n}-\bar \vecv, r}$.  这是译文~\ref{lem:smoothip} (这是译文 $r
  \geq \sqrt{2} \smootheps(q^{-1} \Z^{n})$ 这是译文~$\length{\vecs} \le B$), 这是译文 $\inner{-\vecf,\vecs} + e'$ 这是译文~$6 \eps$  这是译文~$D_{t}$, 这是译文 $t^{2} = r^{2}(\length{\vecs}^{2}+B^{2})$.  这是译文 $e+\inner{-\vecf,\vecs} + e'$ 这是译文~$6 \eps$  这是译文 $D_{(t^2+\alpha^2)^{1/2}}$, 这是译文.
\end{proof}

\begin{comment}

\subsection{Application to Ring-LWE}

\cnote{I think the theorem in this section is a special case of the
  ``ring-LWE embedding'' theorem from GHPS (SCN'12), which says that
  we can embed LWE over a cyclotomic ring with index $m$ into LWE over
  a cyclotomic of index evenly divisible by $m$.  This theorem is the
  special case where $m=1$ (so the cyclotomic ring is just $\Z$).}

\begin{definition}
\label{def:rlwe}
Let~$R_{n,q}$ denote the ring~$\Z_q[x]/(x^n+1)$, with~$n$ a power
of~$2$ and~$q$ an integer. Let~$\chi$ be a distribution
over~$\mathbb{T}_n = \mathbb{T}[x]/(x^n+1)$, and~$\mathcal{D}$ be a
distribution from~$R_{n,q}$. For~$s \in R_{n,q}$, we define the
distribution~$A_{n,q,\chi,s}^R$ as follows: sample~$a \hookleftarrow
U(R_{n,q})$ and~$e \hookleftarrow \chi$; then return~$(a,
\frac{1}{q}(a\cdot s)+e) \in R_{n,q} \times \mathbb{T}_n$.


The decision \emph{Ring Learning with Errors} problem
$\rlwe_{n,q,\chi}(\mathcal{D})$ is as follows: Let~$\vec{s}$ be
sampled from~$\mathcal{D}$; the goal is to distinguish between
arbitrarily many independent samples from~$A_{n,q,\vec{s},\chi}^R$ and
the same number of independent samples from the uniform distribution
over the support of~$A^R_{n,q,\vec{s},\chi}$.
\end{definition}

The above RLWE problem is a special case of the ring learning with
errors problem that was introduced
in~\cite{DBLP:conf/eurocrypt/LyubashevskyPR10}. Since then, it has
proved useful for designing cryptographic schemes with excellent
asymptotic performance. However, all known hardness proofs for \rlwe
rely on the assumed hardness of standard worst-case lattice problems
restricted to the family of so-called ideal lattices. In our present
setup, an ideal lattice is an ideal of~$\Z[x]/(x^n+1)$.
Theorem~\ref{th:reverse} leads to a reduction from \lwe to \rlwe, and
thus provides a hardness proof for \rlwe under standard worst-case
lattice problems over all lattices.

\begin{theorem}
\label{th:rlwe}
For any~$q, \chi, \mathcal{D}$, there exists a polynomial-time
reduction from~$\lwe_{1,q,\chi}(\mathcal{D})$
to $\rlwe_{n,q,\chi'_n}(\mathcal{D}_n)$ where~$\chi'_n = \chi'+\chi' x
+ \cdots + \chi' x^{n-1}$, $\chi'$ is the sum of $n$ independent
copies of~$\chi$ and~$\mathcal{D}_n = \mathcal{D} + \mathcal{D} x +
\cdots + \mathcal{D} x^{n-1}$.
\end{theorem}

\begin{proof}
  For~$s_0,\ldots,s_{n-1} \in \Z_p$, we define the
  distribution~$A_{n,q,\chi,(s_i)_i}$: Sample~$a \hookleftarrow
  U(\Z_q)$ and~$e_0,\ldots,e_{n-1} \hookleftarrow \chi$, and
  return~$(a, \frac{1}{p}(a\cdot
  s_0)+e_0,\ldots,\frac{1}{p}(as_{n-1})+e_{n-1}) \in \Z_p \times
  \mathbb{T}^n$.  Using a standard hybrid argument, we see
  that~$\lwe_{1,q,\chi}(\mathcal{D})$ reduces to distinguishing between
  arbitrarily many independent samples from~$A_{n,q,\chi,(s_i)_i}$ and
  the same number of independent samples from the uniform distribution
  over the support of~$A_{n,q,\chi,(s_i)_i}$. We describe an efficient
  mapping of~$n$ elements~$(a_i, (b_{ij})_{0\leq j<n})_{0 \leq i<n}$
  from~$\Z_p \times \mathbb{T}^n$ to a single element in~$R_q \times
  \mathbb{T}_n$ such that the $A_{n,q,\chi,(s_i)_i}$ distribution
  (resp.\ uniform distribution over the support
  of~$A_{n,q,\chi,(s_i)_i}$) is mapped to the~$A_{n,q,\chi_n,s}^R$
  distribution (resp.\ uniform distribution over the support
  of~$A_{n,q,\chi_n,s}^R$), where~$s = s_0+s_1x + \cdots
  +s_{n-1}x^{n-1}$. The mapping is as follows:
\begin{itemize}
\item Define~$a = a_0 + a_1x + \cdots + a_{n-1}x^{n-1} \in R_{n,q}$.
\item Define~$b = \sum_{0\leq k<2n-1} (\sum_{i+j = k} b_{ij}) x^k \bmod (x^n+1)$.
\item Return~$(a,b)$.  
\end{itemize}
Clearly, the reduction runs in polynomial-time, and the uniform
distribution is mapped to the uniform distribution. Finally, in case
the input~$b_{ij}$'s are of the form~$\frac{1}{p} (a_i\cdot s_j) +e_j$
for all~$i,j$, then all the coefficients of~$b - \frac{1}{p}(a\cdot s)
\in \mathbb{T}_n$ are sums of~$n$ independent samples from~$\chi$.
\end{proof}

As a corollary of Theorems~\ref{th:reverse} and~\ref{th:rlwe}, there
exists a (classical) polynomial-time reduction from standard
worst-case problems over lattices with polynomial approximation
factors to~$\rlwe_{n,p^n,D_{\sqrt{n}\alpha}}(U(R_{n,p^n}))$.

\end{comment}

\begin{comment}
\begin{verbatim}
Also, we were wondering about a possible relationship with
   http://www.di.ens.fr/~lyubash/papers/subsetsumcrypto.pdf
(digit-wise expansion of a large integer, along columns rather than rows).

Notation: \vec* denotes an n-dimensional vector mod p (i.e., in Zp^n),
while capital letters (A,S) denote scalar values mod p^n (i.e., in
Z_{p^n}).

The basic approach is to map LWE samples (\vecr,b) in Zp^n \times Zp
to (A, b*p^{n-1} + noise) in Z_{p^n} \times Z_{p^n}, where A is chosen
from r using some *randomized* encoding (not deterministically).  Then
the amount by which b*p^{n-1} differs from the value A*S will be a
centered Gaussian, and we only need to add a small amount of extra
noise (or maybe none at all?) to smooth it out.

Details:

For A,S in Z_p^n we can write A = sum_i p^i * Ai for *integers* ai
(not necessarily in {0,...,p-1}), and similarly for S.

Then the product A*S mod p^n can be written as \veca^t \matP \vecs mod
p^n, where

\veca = (a_{n-1}, ..., a_0),   \vecs = (s0, s1, ..., s_{n-1}),  and
\matP is the fixed lower-triangular Toeplitz matrix with
\vecp=(p^{n-1}, p^{n-2}, ..., p^0) down its leftmost column.

Ideally, we would then like to choose \veca (defining A) to satisfy

  \veca^t \matP = p^{n-1} * \vecr^t    (mod p^n)

where \vecr is the "left-hand side" of the LWE challenge.

If we could do this, then the product A*S would be exactly p^{n-1} *
<r,s>.  We are given the noisy value b=<r,s>+error, so we could just
plug in p^{n-1}*b (mod p^n) to complete the reduction.  Notice that
this would keep the error rate exactly the same, and moreover the
error term would be divisible by p^{n-1}, which seems like too much to
ask.... and it is.

The problem, of course, is that it is usually not possible to solve
the above system, because \matP is highly non-invertible mod p^n.  (It
doesn't even help that the RHS is a multiple of p^{n-1}.)

Luckily, it's enough to approximately solve the system, as long as the
"error" \veca^t \matP - p^{n-1} * \vecr^t is a centered Gaussian and
\vecs is short.  And the structure of \matP makes this easy to do --
Oded's formulas are the answer.

[Chris note: the formulas come from an earlier email, with slightly
different notation.  But they are easy to work out: the entries of
\veca are chosen iteratively from discrete Gaussians over the
integers, with appropriate centers determined by the previously chosen
entries of \veca.]

Another way of going about it is to multiply both sides by \matP^-1
**over the integers**, which has a very simple form (exercise for
reader).  Then we end up with a fractional RHS which we can just
randomized-round to get an integral \veca.  Doing it this way causes
the "error" \veca^t \matP - p^{n-1} * \vecr^t to be a non-spherical
but "nearly spherical" Gaussian (because \matP is almost orthogonal),
which would suffice for our purposes.  It's not clear that this
non-adaptive solution has any important advantage here, though.

(Note: these two approaches correspond to randomized versions of
Babai's nearest-plane and "simple rounding" algorithms, respectively.)

\end{verbatim}
\end{comment}



%%% Local Variables: 
%%% mode: latex
%%% TeX-master: "lwehard"
%%% End: 



%
%\begin{sloppypar}
%\paragraph{Acknowledgments} We thank Elette Boyle, Adam Klivans,
%Vadim Lyubashevsky, Sasha Sherstov, and Vinod Vaikuntanathan for
%useful discussions.
%\end{sloppypar}


\bibliographystyle{abbrv}
\bibliography{lattices,crypto,ibe}

%\appendix

%\section{这是译文}

\anote{Still some problems that we don't have time to fix in the proof, so maybe it is better to remove this part unless someone have some time to spend on it (proof of claim~\ref{cl:real_left_size} and in claim~\ref{cl:real_right_side} the distribution of $\eps$ need to be studied conditioned on $\vec{s}$ ?}

\subsection{这是译文}

这是译文 $q$ 这是译文.

这是译文 $\phi$ 这是译文~$\mathbb{T}$ 这是译文~$\vec{s} \in \Z^n$ 这是译文-$A_{\vec{s},\phi}$ 这是译文~$\mathbb{T}^{n+1}$ 这是译文 $\vec{a}$ 这是译文~$\mathbb{T}^{n}$, 这是译文 $e \in \mathbb{T}$ 这是译文~$\phi$, 这是译文~$(\vec{a}, \inner{\vec{a},\vec{s}}+ e \bmod 1)$.

\begin{definition}
  这是译文~$n \geq 1$ 这是译文, $\beta \in (0,1)$ 这是译文~$\Upsilon$ 这是译文~$\mathbb{T}$ (这是译文). 这是译文~$n, \beta$ 这是译文~$\Upsilon$ 这是译文~$\lambda$ 这是译文. 
  
  这是译文 \emph{这是译文}, \emph{这是译文-$\text{LWE}_{\Upsilon}(D_{\Z^n, \beta})$}, 这是译文 $\vec{s} \in \Z^n$ 这是译文 $D_{\Z^n, \beta}$ 这是译文~$\phi$ 这是译文~$\Upsilon$; 这是译文 $A_{\vec{s},\phi}$ 这是译文 $U( \mathbb{T}^{n+1})$.
\end{definition}

这是译文-$\text{LWE}_{\Upsilon}(D_{\Z^n, \beta})$ 这是译文~$\Upsilon$ 这是译文 $\phi$ 这是译文 \emph{这是译文} 这是译文~$\Upsilon$.


\subsection{这是译文$_{n,q,\nu_{\sqrt{2}\alpha}}(D_{\Z^n, \alpha q})$ 这是译文-$\text{LWE}_{n,\Upsilon_{\beta_2}}(D_{\Z^n, \beta_1})$}


\begin{theorem}\label{th:real}  这是译文~$n \geq 1$, $q \in [2, 2^{n^{O(1)}}]$ 这是译文 $\alpha, \beta_1, \beta_2
  \in (0,1)$, 这是译文~$\lambda$.
%  Let~$\mathcal{D}$ be an arbitrary distribution of $\Z_q^n$ and define~$s_{\max}(\mathcal{D})$ 
%  as any real such that~$\Pr_{\vec{s} \hookleftarrow \mathcal{D}}[ \Vert \vec{s} \Vert \geq s_{\max}(\mathcal{D}) ] \leq \lambda^{- \omega(1)}$.
  
  这是译文 $\beta_1 \geq 2 \alpha q $,  $\beta_2 \geq \alpha \mult \omega(\sqrt{n \log \lambda} + \log \lambda)$ 这是译文~$\alpha q \geq \omega(\sqrt{\log \lambda})$.  这是译文 \emph{这是译文$_{n,q,\nu_{\sqrt{2}\alpha}}(D_{\Z^n, \alpha q})$}  这是译文 \emph{这是译文-$\text{LWE}_{n,\nu_{\leq \beta_2}}(D_{\Z^n, \beta_1})$}.
\end{theorem}

\begin{proof}
这是译文-$\text{LWE}_{n,\nu_{\leq \beta_2}}(D_{\Z^n, \beta_1})$ 这是译文 $\lambda^{-O(1)}$. 这是译文 $m$ 这是译文~$U(\Z_q^n \times \mathbb{T})$ 这是译文~$A_{\vec{s},\nu_{ \alpha}}$,  这是译文~$\vec{s}$ 这是译文 $\mathcal{D}= D_{\Z^n,\alpha q}$. 这是译文 $ \Vert \vec{s} \Vert \leq \alpha q \mult \omega(\sqrt{n} + \sqrt{\log \lambda})$, 这是译文~$\mathcal{R}$ 这是译文: 
 
 \begin{itemize}
\item[$\bullet$] 这是译文 $\vec{s}_1$ 这是译文 $D_{\Z^n, \sqrt{\beta_1^2 - \alpha^2 q^2} }$.
\item[$\bullet$] 这是译文 $(\vec{a},b)$, 这是译文 $(\vec{a}',b')$ 这是译文:
\begin{itemize}
\item[$\bullet$] 这是译文 $\sigma = \frac{1}{q} \omega (\sqrt{ \log \lambda })$, 这是译文 $\vec{a}' = \frac{1}{q} \vec{a} + \nu_{\sigma} (\bmod~1)$ 这是译文 $b' = b + \inner{\vec{a}', \vec{s}_1}$.
\item[$\bullet$] 这是译文 $(\vec{a}',b')$.
\end{itemize}
\item[$\bullet$] 这是译文-$\text{LWE}_{n,\nu_{\leq \beta_2}}(D_{\Z^n, \beta_1})$ 这是译文.
\end{itemize} 这是译文~$b \in \mathbb{T}$ 这是译文~$\frac{1}{q}
\inner{\vec{a},\vec{s}}+ e$ 这是译文~$e \hookleftarrow \nu_{\alpha}$. 这是译文: \begin{equation}
\label{eq:real_rewrite_b}
\left(\vec{a'}, b\right) = 
\left(
\vec{a'}, \inner{\vec{a}',\vec{s} + \vec{s}_1} +  \inner{\vec{\bfeps},\vec{s}} + e  \right).  
\end{equation} 这是译文~$(\vec{a}',\vec{\bfeps})$.  

\begin{claim}
\label{cl:real_left_size} 这是译文~$\sigma \geq \frac{1}{q}\omega(\sqrt{\log \lambda})$. 这是译文~$\vec{a}'$ 这是译文~$U(\mathbb{T}^n)$, 这是译文~$\overline{\vec{a}'} \in \mathbb{T}^n$, 这是译文~$\vec{\bfeps}$ 这是译文~$\vec{a}' =
  \overline{\vec{a}'}$ 这是译文~$D_{ \overline{\vec{a}'} + \frac{1}{q} \Z^n,\sigma}$.
\end{claim}


\begin{proof}
TODO
\end{proof}

这是译文~$\vec{s} + \vec{s}_1$ 这是译文~$\inner{\vec{\bfeps},\vec{s}} + e $ 这是译文.~\eqref{eq:real_rewrite_b}, 这是译文~$\vec{a}'$.

\begin{claim}
\label{cl:real_right_side}  这是译文~$\beta_1 \geq 2 \alpha q$ 这是译文~$\alpha q \geq \omega ( \sqrt{\log \lambda})$, 这是译文~$\vec{s} + \vec{s}_1$ 这是译文~$D_{\mathbb{Z}^n,\beta_1} $.

这是译文~$\vec{\bfeps} \hookleftarrow \nu_{\sigma}$ 这是译文~$e \hookleftarrow \nu_{\alpha}$, 这是译文~$ \inner{\vec{\bfeps},\vec{s}} + e$ 这是译文~$\nu_{\beta}$ 这是译文~$\beta=\sqrt{\sigma^2\|\vec{s}\|^2+ \alpha^2}$.

\end{claim}

\begin{proof}
 这是译文~$\beta_1 \geq 2 \alpha q$ 这是译文~$\alpha q \geq \omega ( \sqrt{\log \lambda})$, 这是译文~$\sqrt{\beta_1^2 - \alpha^2 q^2} \geq \alpha q \geq \omega ( \sqrt{\log \lambda})$. 这是译文~\ref{le:discrete_plus_cont_Gauss} 这是译文~$\vec{s} + \vec{s}_1$ 这是译文~$D_{\mathbb{Z}^n,\beta_1} $.


这是译文~$\vec{\bfeps} \hookleftarrow \nu_{\sigma}$ , 这是译文 $\inner{\vec{\bfeps},\vec{s}}$ 这是译文  $\nu_{\sigma\|\vec{s}\|}$. 这是译文 $ \inner{\vec{\bfeps},\vec{s}} + e$ 这是译文~$\nu_{\beta}$ 这是译文~$\beta=\sqrt{\sigma^2\|\vec{s}\|^2+ \alpha^2}$.
\end{proof}

这是译文~$\vec{s}' = \vec{s} + \vec{s}_1$. 这是译文~$\vec{a}'$ 这是译文~$U(\mathbb{T}^n)$, 这是译文~$b' = \inner{\vec{a}',\vec{s}'} + e'$, 这是译文~$e'$ 这是译文~$\vec{a}'$ 这是译文~$\nu_{\sqrt{\sigma^2\|\vec{s}\|^2+ \alpha^2}}$. 这是译文~$\sigma = \frac{1}{q} \omega (\sqrt{\log \lambda})$ 这是译文~$ \Vert \vec{s} \Vert \leq \alpha q \mult \omega(\sqrt{n} + \sqrt{\log \lambda})$. 这是译文, $\sqrt{\sigma^2\|\vec{s}\|^2+ \alpha^2} \leq \alpha \mult \omega(\sqrt{n \log \lambda} + \log \lambda)$.

这是译文~$\vec{s}'$ 这是译文~$D_{\mathbb{Z}^n,\beta_1}$, 这是译文~$\beta_2$. 这是译文~\ref{th:real}.


\end{proof}



\end{document}

%%% Local Variables:
%%% mode: latex
%%% TeX-master: t
%%% End:
